% 大连理工大学数学专业科研修订稿：现行官方模板附件尚需核对。
\newif\ifdraftfallback
\IfFileExists{ctexbook.cls}{
  \documentclass[UTF8,12pt,oneside]{ctexbook}
}{
  \draftfallbacktrue
  \documentclass[12pt,oneside]{book}
  \usepackage{fontspec}
  \setmainfont{Noto Serif CJK SC}[AutoFakeSlant=.15]
  \setsansfont{Noto Sans CJK SC}
  \setmonofont{DejaVu Sans Mono}[Scale=.85]
  \XeTeXlinebreaklocale "zh"
  \XeTeXlinebreakskip=0pt plus 1pt
  \renewcommand{\contentsname}{目录}
  \renewcommand{\figurename}{图}
  \renewcommand{\tablename}{表}
  \renewcommand{\bibname}{参考文献}
}
\usepackage[a4paper,margin=2.5cm]{geometry}
\usepackage{amsmath,amssymb,amsthm,graphicx,booktabs,longtable}
\renewcommand{\proofname}{证明}
\usepackage[hidelinks]{hyperref}
\ifdraftfallback
  \usepackage{titlesec}
  \titleformat{\chapter}[hang]{\centering\sffamily\bfseries\Large}{\thechapter}{1em}{}
  \titlespacing*{\chapter}{0pt}{0pt}{24pt}
  \titleformat{\section}{\sffamily\bfseries\large}{\thesection}{1em}{}
  \titleformat{\subsection}{\sffamily\bfseries\normalsize}{\thesubsection}{1em}{}
\fi
\hypersetup{pdftitle={局部随机神经网络与自适应采样：本科毕业论文修订稿v3},pdfauthor={作者信息待确认}}

\theoremstyle{definition}
\newtheorem{proposition}{命题}[chapter]
\newtheorem{remark}{注}[chapter]
\newcommand{\R}{\mathbb R}
\newcommand{\norm}[1]{\left\lVert#1\right\rVert}
\newcommand{\figurefile}[3]{\begin{figure}[htbp]\centering\includegraphics[width=.98\textwidth]{figures/#1.pdf}\caption{#2}\label{#3}\end{figure}}
\setlength{\parindent}{2em}
\setlength{\parskip}{0pt}
\linespread{1.25}
\emergencystretch=3em
\pagestyle{plain}
\title{基于自适应采样的局部随机神经网络求解椭圆型界面问题的数值研究\\[1em]\large 本科毕业论文科研初稿（修订版v3）}
\author{作者、学号、学院及指导教师信息待填写}
\date{2026年10月16日}
\begin{document}
\frontmatter
\begin{titlepage}
\centering
{\sffamily\bfseries\LARGE 大连理工大学\par}
\vspace{1.0cm}
{\sffamily\Large 本科毕业论文\par}
\vspace{1.2cm}
{\sffamily\bfseries\Large 基于自适应采样的局部随机神经网络\par
求解椭圆型界面问题的数值研究\par}
\vspace{0.7cm}
{\large Numerical Study of Local Randomized Neural Networks\par
with Adaptive Sampling for Elliptic Interface Problems\par}
\vspace{1.2cm}
\begin{tabular}{ll}
学院：& 数学科学学院（信息待本人确认）\\[.5em]
专业：& 数学与应用数学（信息待本人确认）\\[.5em]
学生姓名：& 待填写\\[.5em]
学号：& 待填写\\[.5em]
指导教师：& 待填写\\[.5em]
完成日期：& 待按实际提交日期填写
\end{tabular}
\vfill
{科研修订稿v3，2026年10月16日\par}
\small 本页为审阅封面，非已核准的学校官方模板。\par
\end{titlepage}
\chapter*{修订稿说明}
本稿使用真实程序输出及公开可核实文献，保留未达到预期优势的实验。学校本科论文规范已开展核查，但本届数学学院细则及现行官方模板附件尚未完整取得，具体差距另存项目文档。封面信息、声明与签名须由本人和指导教师按实际情况核实，不代填指导记录、评语或签名。本页供审阅使用，正式提交时须按学院要求处理。

\chapter*{摘要}
椭圆型界面问题具有分片系数及界面耦合条件，配置点在区域和界面上的分配会影响强形式离散系统的精度与稳定性。本文在既有局部随机神经网络思想基础上，独立实现固定隐藏参数的单隐层双子域方法，通过线性最小二乘确定输出系数。采用解析导数建立内部方程、外边界条件以及界面值和法向通量跳跃条件，并补充有限差分导数模式，对公开文献中的花形界面算例开展近似复现。

首先以正弦制造解验证Poisson方程五点差分和随机特征程序，差分观测收敛阶接近二阶。随后比较固定均匀点、固定随机点、仅PDE残差驱动采样及区域与界面联合指标采样，并开展预算、参数、消融和独立局部峰实验。所有随机比较至少采用五个固定种子，区分物理点预算与标量方程预算，在独立测试点报告解误差、界面误差及完整计算成本。

实验表明，在本文已测试的平滑直线界面问题与预设参数下，联合策略未取得相对均匀、随机基线的稳定解误差优势；增加界面资源的固定随机方案也能解释部分消融收益。局部峰问题中联合策略在部分对照改善，但多数情况下仍未优于均匀基线，且部分窄峰精度不足。尺度平衡、特征宽度、随机权重范围和SVD截断会显著改变误差及有效秩，解误差与物理通量误差可能具有不同排序。固定矩阵扰动实验进一步表明，系数敏感性与预测函数变化不可混为一谈。本文的贡献定位为已有方法的独立实现、可重复比较与有限范围稳定性分析，不主张提出具有普遍优势的原创采样算法。

\noindent\textbf{关键词：}椭圆型界面问题；局部随机神经网络；配置点；残差采样；最小二乘
\chapter*{Abstract}
Elliptic interface problems involve discontinuous diffusion coefficients and coupled solution and flux conditions across an internal interface. The distribution of collocation points can affect both the approximation obtained from these conditions and the conditioning of the resulting linear system. This thesis studies that distribution within an independently implemented local randomized neural network method. Each subdomain uses fixed hidden parameters and trainable linear output coefficients, which are determined by weighted least squares on a CPU. Analytic derivatives are used in the main experiments, and a finite difference derivative mode is implemented for an approximate reproduction of a published example.

Manufactured solutions provide consistent source terms, exterior boundary data, solution jumps, and coefficient-weighted normal flux jumps. A sine solution on the unit square first verifies the Poisson finite difference baseline and the random feature implementation. The finite difference results exhibit an observed order close to two. Subsequent experiments use a straight interface with continuous or nonzero jump data, diffusion coefficient contrasts in both directions, and additional localized peaks. A flower-shaped interface is included as a separate reproduction study, with implementation assumptions and discrepancies explicitly retained.

Four strategies are compared: fixed uniform points, fixed random points, interior PDE residual sampling, and joint allocation based on interior and interface residual indicators. Hidden features and random seeds are shared within each comparison. Physical point budgets and scalar equation budgets are distinguished because an interface location supplies two scalar conditions. Independent evaluation points are excluded from point selection. Random comparisons use at least five fixed seeds, while the main flower reproduction uses ten. The reported quantities include relative solution error, maximum absolute error, independent PDE residuals, interface errors, numerical rank, conditioning, and the complete adaptive computation time.

For the smooth straight-interface configurations tested here, joint sampling does not consistently outperform the static baselines in solution error. A fixed random control with the same final group allocation also accounts for part of the apparent benefit of an ablation. Localized peak problems produce mixed comparisons, and some narrow peaks remain inaccurate under the chosen fixed feature spaces and budgets. Scaling, feature count, initialization range, and singular value truncation introduce trade-offs between solution accuracy, physical flux accuracy, and numerical conditioning. Fixed-matrix perturbations show that coefficient sensitivity and prediction sensitivity must be distinguished; these observations do not establish general robustness to matrix perturbations.

The mathematical discussion provides self-contained finite-dimensional error identities, conditional stability bounds, and a restricted continuous energy estimate with explicit assumptions. These results do not turn finite training residuals into a general PDE error guarantee. The contribution is a reproducible implementation and a systematic numerical analysis of existing techniques, rather than an assertion of an original sampling algorithm with universal advantages. Further validation requires independent models, fuller access to the original implementations, and confirmation of the current institutional submission requirements.

\noindent\textbf{Keywords:} elliptic interface problem; randomized neural network; collocation; adaptive sampling; least squares
\tableofcontents
\mainmatter
\chapter{绪论}
\section{研究背景与问题}
椭圆型偏微分方程常用于描述稳态扩散和势场。对于由不同材料组成的区域，扩散系数在材料界面发生跳跃，真解通常仅在每个子区域内光滑。即使解连续，梯度也可能不连续；在给定非零跳跃数据时，解本身亦可不连续。因此，单纯在整个区域内使用一个光滑近似并忽略界面条件，不能完整表达所要求解的问题。

传统差分或有限元方法通过网格及界面处理建立离散系统。随机神经网络方法则预先生成隐藏层参数，把近似写为固定特征的线性组合，只求输出系数。对于线性PDE，内部方程和线性边界、界面条件均可转化为线性约束。这为CPU上的可重复小规模研究提供了条件，但隐藏特征的相关性、配置点覆盖及约束尺度会影响求解结果。

本文研究的问题是：在固定随机特征和明确计算预算下，如何评价内部与界面配置点的选取及分配？残差较大的位置增加点数是否必然改善解误差？当系数对比度提高时，误差、矩阵条件性和计算成本是否具有一致的变化趋势？这些问题不能仅用训练损失回答，需要独立真解验证与合理静态基线。

\section{相关研究及贡献边界}
Sun、Dong和Wang将局部随机网络与间断Galerkin框架结合，提出多个格式并给出带假设的误差分析\cite{dgpreprint}。其中弱形式积分、面耦合及近似条件与本文强形式配置不同，不能将两者统称为同一离散格式。Li和Wang使用不同子域的随机网络求解界面问题，采用中心差分近似特征导数，并在随机点上叠置内部、界面和边界约束\cite{interfacepreprint}。其预印本总结明确把自适应配置点、初始化范围及数值分析列为后续方向；这为本文研究动机提供依据，却不保证本文方案在现阶段具有新颖性。

Dang、Wang和Jiang的AG-RaNN使用频率信息、残差和增长策略构造新增特征，并涉及域分裂\cite{ag}。本文阅读的是2025年的第三版预印本。其特征或网络架构适应与本文固定特征下的空间点适应不同；本文不直接套用其在特定正则性、参数及采样假设下的理论结论。Shang等结合单位分解、PCA与重叠Schwarz预条件来改善局部随机特征系统求解\cite{schwarzpreprint}，提示条件性不仅与点选取有关，也与特征冗余和求解器有关。

残差驱动采样本身已有明确先行研究。Wu等比较固定采样及RAD、RAR-D等方法\cite{rar}。较近的Residual-Christoffel方法依据算子作用后的特征空间几何构造密度、权重和白化\cite{christoffel}，其中“残差空间”并不等于当前近似解残差的幅值。本文的简单联合指标与这些技术的差异尚不足以支持原创算法声明。RNN-PG 2023论文正式书目信息已核实\cite{pg}，但全文尚未获得，因此本文不依据其未读推导建立算法或归纳作者结论。

固定隐藏参数、线性确定输出参数的思想具有更早的研究基础。Huang、Zhu和Siew的极限学习机论文采用随机隐藏节点及解析输出权重\cite{elm2006}；本文仅根据已核实的出版信息与摘要引用这一基本定位，不转述未阅读全文中的具体定理。Anderson等的ELM-FBPINNs结合局部随机特征、单位分解及多层结构\cite{elmfbp}，也提醒我们把空间局部化、特征扩充和选点三个机制分别考察。本文已读公开第三版，与早期版本的结果不混用。

\section{研究内容及组织}
本文的工作包括四个部分。第一，独立实现Poisson五点差分和随机特征配置，验证数学数据与导数。第二，构造零与非零跳跃的双子域制造解，研究加权最小二乘、尺度平衡与特征参数。第三，在共享特征和明确预算下比较四种采样及消融，保留无效、恶化和局部峰精度不足的结果。第四，给出标准有限维稳定性结论的自含推导，并通过扰动及花形界面近似复现讨论数值限制。

第二章介绍模型、制造解、差分和最小二乘；第三章推导局部随机特征离散系统；第四章说明采样及预算机制；第五章报告真实数值实验；第六章总结有限范围结论与未解决问题。附录给出补充模型与复现路径。

\section{研究范围与学术真实性}
本文不要求原创定理或重大算法创新。实验均为程序实际输出，原文报告数值与本项目数据分别列示。所有比较至少五种子，代表图不依据最好误差挑选。正式版和预印本存在题名、年份差异时分别记录；无法核实的全文内容不补写。本文尚未开展真实外部数据、三维、多层增长或完整Schwarz实现，也不作优于PINN或传统界面有限元的直接数值结论。

\section{研究问题的细化及评价原则}
\subsection{为何把配置点选择单独作为研究对象}
在局部随机特征方法中，隐藏层参数、子域划分、配置点、损失权重和最小二乘求解器都能影响精度。如果同时调整所有因素，即使最终误差降低，也难以判断收益来自更好的表达空间、更多约束，还是更合适的选点。本文因此固定两个子域及隐藏特征，再考察点的数量和位置，把问题缩小为本科阶段可以独立验证的有限维研究。

这个缩小并没有消除所有交互。局部特征在不同位置可能具有不同的导数尺度，非均匀点又会改变观测矩阵，截断SVD随之改变有效空间。因此所谓固定网络并不意味着数值上保留的方向完全固定。本文报告奇异值、数值秩和物理误差，目的就是避免把所有变化都归因于残差指标。

界面问题还多了一种资源竞争：内部点主要约束各侧PDE，界面点同时约束值和通量。给界面更多物理位置会产生更多标量方程，而给内部更多点可能改善区域覆盖。对一种资源有利的安排，可能减少另一种资源。因此本课题的核心不是简单地证明残差越大越应多放点，而是检验这种直觉在预算、尺度与固定空间共同限制下是否成立。

\subsection{三个可检验的问题}
第一个问题是正确性：从模型、导数到矩阵的全部环节是否一致？这一问题由制造解一致性、非零跳跃、差分收敛和单元测试回答。没有通过这一层的算法比较没有解释基础；特别是高对比度时，小解误差也不能替代通量行的单独检验。

第二个问题是效果：在相同最终资源下，联合指标是否比静态点和仅PDE指标更有利？本文同时采用点预算与方程预算，并在每个设置内共享隐藏特征和随机种子。比较指标不仅包括相对解误差，还包括最大误差、独立PDE残差、两类界面误差及时间。不同指标不强制给出相同排序。

第三个问题是解释：如果联合方案无效或只有局部收益，哪些观测可以排除过强的解释？同配额固定随机对照可检验残差选点是否为消融收益的必要条件，宽度与范围扫描揭示表达空间和条件性变化，局部峰避免只在平滑低频模型上评价。它们仍不能构成完整因果证明，但能限制结论的范围。

\subsection{研究贡献的三个层次}
第一层是实现贡献，即把已有数学方法独立组织成可运行、可追溯的程序，并用明确模型验证。这一层要求代码与文稿一致，不能因为技术已有就省略复现细节，也不能把实现已有技术说成发明该技术。

第二层是实验与分析贡献，即在明示预算和参数下获得可重复的比较，解释观察到的取舍，保留与预期不一致的结果。对于本科论文，这种工作可以体现综合运用分析、代数、数值计算与概率统计知识的能力。其价值在于回答具体问题，而不是实验次数越多就自动意味着结论越强。

第三层才是方法改进贡献。只有相对于合理基线、在独立问题上重复验证效果，并排除额外资源或参数调优等解释后，才适合提出有限范围的改进主张。本文原计划探索这一层，但当前结果尚未达到相应证据门槛，因此最终定位以第二层为主。这是根据真实结果调整论断，不是放弃原来的研究问题。

\section{制造解研究的适用范围}
制造解允许精确计算真解、源项和跳跃，从而区分代码错误与近似误差。本文不依赖真实材料观测，也不需要训练数据标签；PDE及界面数据本身构成求解约束，解析真解只用于构造测试问题和独立误差评估，不进入残差选点。这样的设置适合验证数值算法，但与未知真解的真实应用仍有距离。

制造解的选择也会形成偏差。如果只使用非常平滑且与随机特征容易匹配的函数，配置点差异可能很小；如果使用超过当前固定空间分辨能力的尖峰，所有选点策略都可能精度不足。本文因此同时保留平滑、局部峰和复杂界面三类问题，而不把其中一类的排序外推到全部椭圆型问题。

系数对比度的方向需要区分。把$\beta=(1000,1)$换成$(1,1000)$不仅改变某个比值的写法，还改变高系数所在侧及制造解中系数依赖项。当前模型虽有规则几何，真解并非在所有参数下严格左右对称，因此正反向都实际运行。这比只报告单一方向更有利于检查尺度和组分配的影响。

需要注意，具有解析解不等于连续误差可以被机器精确积分。本文使用独立规则点或Gauss型积分近似范数，另做密度敏感性检查。图表中的数值是明确验证规则下的结果，不是经过区间算术认证的全域上下界。对于极小误差及高度振荡界面，积分规则和浮点求值本身也可能成为结果解释的一部分。

本论文最后形成的是可供逐项复核的计算研究。数学证明只覆盖已写明假设，数值结论只覆盖真实运行，尚未取得的原文和本届院系细则列为未决事项。这样的边界使读者能够知道哪些内容可以直接复算，哪些仍需要更强理论、作者代码或新的独立测试。


\chapter{预备知识与数学模型}
\section{椭圆型界面问题}
设单位正方形$\Omega=(0,1)^2$被$\Gamma=\{x=1/2,0<y<1\}$分为$\Omega_1$与$\Omega_2$，法向$n=(1,0)$从左指向右。系数$\beta_i>0$在子域内为常数。本文约定跳跃$[v]=v_2-v_1$，与参考预印本的相反约定须同时转换约束行和数据。
\begin{equation}\label{eq:interface}
\begin{cases}
-\beta_i\Delta u_i=f_i,&\text{在 }\Omega_i,\ i=1,2,\\
u=g_D,&\text{在 }\partial\Omega,\\
[u]=g_u,\quad[\beta\nabla u\cdot n]=g_q,&\text{在 }\Gamma.
\end{cases}
\end{equation}
本研究使用分片光滑制造解，把各项数据从同一函数严格导出。方程在每侧求导，不对跨界面不连续的总函数直接应用普通中心差分。

\section{符号与分析层次}
表\ref{tab:symbols}统一本文的主要记号。内部、外边界和界面是不同采样对象；$\Gamma$上的一个物理位置会产生两条标量约束。所有奇异值和条件数均针对当前加权矩阵，而非未注明权重的抽象PDE算子。
\begin{table}[htbp]\centering\small
\caption{主要符号及含义}\label{tab:symbols}
\begin{tabular}{ll}\toprule
符号 & 含义\\\midrule
$\Omega_i,\Gamma,n$ & 子区域、界面、从第1侧指向第2侧的单位法向\\
$\beta_i,f_i,g_D,g_u,g_q$ & 扩散系数、源项、边界及两类跳跃数据\\
$N_i,K=N_1+N_2$ & 每侧及总随机特征数\\
$m_k,m_{\rm eq}$ & 第$k$组点数、叠置系统标量行数\\
$\lambda_k,s_k$ & 目标中的平方权重、固定残差尺度\\
$B,d,c$ & 加权矩阵、加权数据、输出系数\\
$\tau,\Pi_\tau$ & 相对SVD阈值、保留右奇异子空间的正交投影\\
$P,v$ & 独立验证点预测矩阵、对应精确解值向量\\
$\rho,\rho_0$ & 连续抽样密度与目标密度（第四章简化情形）\\
$q(x),p(x)$ & 制造解中的多项式；不同于奇异向量$q_j$\\
\bottomrule\end{tabular}\end{table}

需要区分三个分析层次。连续模型规定式\eqref{eq:interface}；有限维离散把它变为配置点约束；浮点求解产生实际系数。本文验证了具体模型数据和离散组装，并分析固定矩阵的一些性质，尚未建立从有限配置点残差到一般连续解范数的统一误差定理。将这三个层次分开，才能避免把线性求解准确误写为PDE整体准确。

\section{Poisson制造解与差分验证}
令$u=\sin(\pi x)\sin(\pi y)$。由于$u_{xx}=u_{yy}=-\pi^2u$，有$f=2\pi^2u$及零Dirichlet边界。每向$n$个内部节点、$h=1/(n+1)$时五点格式为
\begin{equation}\label{eq:fdm}
\frac{4U_{ij}-U_{i-1,j}-U_{i+1,j}-U_{i,j-1}-U_{i,j+1}}{h^2}=f_{ij}.
\end{equation}
按$x$方向快序铺平，矩阵为$A=I\otimes T+T\otimes I$，其中$T=h^{-2}\operatorname{tridiag}(-1,2,-1)$。零边界不增加右端，非零边界则须把已知邻点贡献移至右端。

\begin{proposition}[正弦算例的离散误差]
对上述节点和制造解，五点差分解恰为
\[U=(2\pi^2/\lambda_h)u|_{\rm grid},\qquad\lambda_h=8h^{-2}\sin^2(\pi h/2).\]
相对离散$\ell^2$误差为$\pi^2h^2/12+O(h^4)$。
\end{proposition}
\begin{proof}
一维节点正弦是$T$的特征向量，特征值$4h^{-2}\sin^2(\pi h/2)$；二维张量积对应两个特征值之和。右端为$2\pi^2u|_{\rm grid}$，故解的倍数关系成立。由正弦平方的Taylor展开得
\[\lambda_h=2\pi^2-\pi^4h^2/6+O(h^4),\]
展开$2\pi^2/\lambda_h-1$即得结论。
\end{proof}
这是经典差分性质在特定真解上的初等推导，不是一般PDE或随机网络的新收敛定理。

\section{直线界面的分片制造解}
记$q=x(1-x)$、$p=q(x)(x-1/2)$，定义
\begin{align}
u_i&=\left(\sin\pi x+a_iq+\frac{t_i}{\beta_i}p\right)\sin\pi y,\label{eq:manufactured}\\
f_i&=\left[2\beta_i\pi^2\sin\pi x+\beta_i a_i(2+\pi^2q)-t_i(3-6x)+t_i\pi^2p\right]\sin\pi y.
\end{align}
因为$q''=-2$、$p''=3-6x$，式\eqref{eq:manufactured}满足内部方程。外边界为零，界面上$p=0,q=1/4,p'=1/4$，于是
\begin{equation}
g_u=\frac{a_2-a_1}{4}\sin\pi y,\qquad g_q=\frac{t_2-t_1}{4}\sin\pi y.
\end{equation}
连续模型取$a_1=a_2=0,t_1=t_2=1$，此时通量本身为$\sin\pi y/4$，并非恒零，能有效检验通量行。非零跳跃取$a_1=0,a_2=.4,t_1=1,t_2=1.2$，所以$g_u=.1\sin\pi y,g_q=.05\sin\pi y$。

\figurefile{model_geometry}{直线界面、法向及非零跳跃制造解剖面；此图为解析模型示意，不是网络数值解。}{fig:model_geometry}

\section{固定特征最小二乘与截断}
最小二乘、SVD和伪逆属于标准数值线性代数内容\cite{golub}。以下给出本文所需结论的自含推导，不将其列为原创定理。对实矩阵$B\in\R^{m\times n}$，求$\min_c\norm{Bc-d}_2^2$。SVD写为$B=Q\Sigma V^T$，保留$\sigma_j>\tau\sigma_1$的方向后，截断解为
\begin{equation}\label{eq:tsvd}
c_\tau=\sum_{\sigma_j>\tau\sigma_1}\frac{q_j^Td}{\sigma_j}v_j.
\end{equation}
满列秩时$\kappa_2(B)=\sigma_1/\sigma_n$。由正规矩阵特征值$\sigma_j^2$可知$\kappa_2(B^TB)=\kappa_2(B)^2$，这解释了为什么不直接用正规方程替代QR/SVD。本文使用SciPy的gelsd型求解，并单独记录奇异值诊断\cite{scipy}。

\begin{proposition}[固定矩阵右端扰动]
若$B$满列秩、$c=B^\dagger d$、$\widetilde c=B^\dagger(d+\delta d)$，则$\norm{\widetilde c-c}_2\le\norm{\delta d}_2/\sigma_n$。固定截断子空间时以最小保留奇异值替换$\sigma_n$。
\end{proposition}
\begin{proof}
差值为$B^\dagger\delta d$。由SVD，$\norm{B^\dagger}_2=1/\sigma_n$，应用算子范数不等式即得。固定截断空间的线性映射具有相同推导。若扰动沿最小保留左奇异向量，上界在精确算术中取到。
\end{proof}
这一标准线性代数结论不覆盖矩阵扰动或秩变化。若预测矩阵为$P$，函数变化是$P\delta c$，所以大的系数变化也不自动意味着大的函数变化。岭目标$\norm{Bc-d}^2+\alpha\norm c^2$对奇异方向使用$\sigma/(\sigma^2+\alpha)$，抑制放大但引入偏差，不能保证更小真解误差。

\section{有限维误差分解及其限制}
下面只使用固定矩阵和标准SVD恒等式\cite{golub}，为实验提供可检查的解释。设$B\in\R^{m\times K}$，保留集合非空，$B_\tau^\dagger$为固定截断伪逆，$\Pi_\tau$为对应右奇异向量张成空间上的正交投影，因而$B_\tau^\dagger B=\Pi_\tau$。验证矩阵$P\in\R^{n_v\times K}$作用于系数得到独立验证点的解值，精确值为$v$。

\begin{proposition}[固定离散系统的误差恒等式]\label{prop:finite_error}
给定任意参考系数$c_{\rm ref}$，记$r_{\rm ref}=d-Bc_{\rm ref}$。数据扰动为$\delta d$，实际系数写为
\[\widehat c=B_\tau^\dagger(d+\delta d)+\epsilon_c.\]
则有
\begin{align}
P\widehat c-v={}&(Pc_{\rm ref}-v)-P(I-\Pi_\tau)c_{\rm ref}\nonumber\\
&+PB_\tau^\dagger(r_{\rm ref}+\delta d)+P\epsilon_c.\label{eq:error_decomposition}
\end{align}
从而
\begin{align}
\norm{P\widehat c-v}_2\le{}&\norm{Pc_{\rm ref}-v}_2+\norm{P(I-\Pi_\tau)c_{\rm ref}}_2\nonumber\\
&+\norm{PB_\tau^\dagger}_2(\norm{r_{\rm ref}}_2+\norm{\delta d}_2)
+\norm P_2\norm{\epsilon_c}_2.\label{eq:finite_bound}
\end{align}
\end{proposition}
\begin{proof}
代入$d=Bc_{\rm ref}+r_{\rm ref}$，利用$B_\tau^\dagger B=\Pi_\tau$，得到$\widehat c=\Pi_\tau c_{\rm ref}+B_\tau^\dagger(r_{\rm ref}+\delta d)+\epsilon_c$。左右乘$P$并加减$Pc_{\rm ref}$即得恒等式；三角不等式和算子范数性质给出式\eqref{eq:finite_bound}。
\end{proof}

第一项衡量参考特征组合在验证点上的近似，第二项是截断丢弃分量，第三项包含参考组合的离散残差与数据扰动，第四项描述系数求解误差的传播。这个分解不要求真解属于特征空间，但选取$c_{\rm ref}$及控制各项本身并不自动容易。它也没有分析矩阵装配误差、候选点选择的概率保证或连续Sobolev范数。本文将其作为标准有限维代数的自含整理，不列为原创收敛定理。

特别地，真实函数值近似好，不保证其二阶导数近似也好；参考组合的$\norm{r_{\rm ref}}$仍需考虑。反之，若只在训练点测到小残差，也不能据此控制未知的$P$或连续全域误差。式\eqref{eq:error_decomposition}说明为什么表达力、配置点、截断和浮点计算必须分别检查。

\section{新增约束与条件数的非单调性}
\begin{proposition}[固定列空间的增行性质]\label{prop:rows}
设旧矩阵$B$与新增行$C$具有相同列数，$\widetilde B=[B;C]$，旧行权重保持不变。以$\min_{\norm z=1}\norm{Bz}$定义最小奇异值，必要时包含零，则$\sigma_{\min}(\widetilde B)\ge\sigma_{\min}(B)$。但条件数不必降低。
\end{proposition}
\begin{proof}
对每个单位向量$z$，$\norm{\widetilde Bz}^2=\norm{Bz}^2+\norm{Cz}^2\ge\norm{Bz}^2$，取最小值得第一结论。取$B=I_2,C=(L,0)$，增行后奇异值为$\sqrt{1+L^2}$和1，条件数由1变为$\sqrt{1+L^2}$，因此并无降低保证。
\end{proof}

这一定性例子说明，更多行既可能提高最小奇异值，也可能提高最大奇异值，不能只看点数推断条件数。更重要的是，本算法按组均方归一化：增加组内点数$m_k$会把旧行的因子从$m_k^{-1/2}$改为新的更小因子。新系统不只是保持旧行不变的简单增行，因此第一结论也不能直接用于当前自适应矩阵。相对截断$\tau\sigma_1$还随最大奇异值变化，有效秩的变化须实际记录。

\section{模型的唯一性及连续残差能说明什么}
\subsection{界面法向与分部积分}
界面条件同时连接两个区域，法向约定必须在模型、离散和误差计算中一致。设$n$从$\Omega_1$指向$\Omega_2$，则两个子域在界面上的外法向分别为$n$和$-n$。对于连续试验函数$v$，两侧分部积分产生的界面项为
\[
\int_\Gamma \beta_1\partial_n u_1v-\int_\Gamma\beta_2\partial_n u_2v
=-\int_\Gamma g_qv.
\]
因此，若解连续、外边界为零，经典形式对应
\begin{equation}\label{eq:weak_interface}
\sum_{i=1}^2\int_{\Omega_i}\beta_i\nabla u_i\cdot\nabla v
=\sum_{i=1}^2\int_{\Omega_i}f_iv-\int_\Gamma g_qv.
\end{equation}
这个负号来自本文的跳跃定义，而非通量条件本身具有另一种物理含义。若改用左减右的跳跃约定，右端符号随之改变。强形式配置不需要实际计算这些积分，但分部积分提供了独立检查通量符号的途径。

非零值跳跃时，解不属于通常的全域连续空间，不能直接把$u$当作齐次Dirichlet能量空间中的函数。本文使用两侧独立特征，正是为了保留这个分片性质。弱形式仍可通过适当的跳跃提升处理；不同弱离散如何惩罚值跳跃属于另一组方法选择，不等同于本文的强配置约束。

\begin{proposition}[直线界面经典解的唯一性]\label{prop:unique}
假设两组解在每侧闭区域上二阶连续，满足相同的内部、外边界、值跳跃和通量跳跃数据。扩散系数为正常数。则两组解相同。
\end{proposition}
\begin{proof}
令差为$w_i$。它满足$-\beta_i\Delta w_i=0$、外边界为零、$[w]=0$及$[\beta\partial_nw]=0$。两侧分别乘$w_i$并分部积分；外边界项为零，界面上$w_1=w_2$，两侧通量项之和也为零。故
\[
\sum_i\beta_i\int_{\Omega_i}|\nabla w_i|^2=0.
\]
由$\beta_i>0$知每侧梯度为零，因而每侧为常数。连续界面使常数相同，零外边界使常数为零。
\end{proof}
该命题只证明在给定光滑类别中的唯一性，没有单独证明任意粗糙数据下的存在性。本文通过先指定分片光滑函数再制造数据，已经给出了所用算例的存在性，因此无需用更强的一般存在理论替代实际模型核查。这个证明也说明，漏掉值条件或通量条件会改变问题，而不能仅视为少加一些训练行。

\subsection{带精确边界和值条件的条件性残差界}
为了区分连续残差与有限训练损失，考虑一个比实际计算更理想的情况。设真解$u$与近似$v$均分片光滑，并且$v$精确满足外边界和值跳跃。记$e=v-u$，连续内部残差$r_i=-\beta_i\Delta v_i-f_i$，连续通量残差$s=[\beta\partial_nv]-g_q$。这里连续范数使用区域面积及界面长度，而不是训练点的均方。

由于$e$在界面连续且在外边界为零，对固定$y$可沿$x$积分：$e(x,y)=\int_0^x e_x(t,y)\,dt$，积分经过界面时由连续性连接两段。Cauchy--Schwarz不等式给出
\[
\norm e_{L^2(\Omega)}\le\norm{\nabla e}_{L^2(\Omega)},\qquad
\norm{e|_\Gamma}_{L^2(\Gamma)}\le 2^{-1/2}\norm{\nabla e}_{L^2(\Omega)}.
\]
这些常数是为当前正方形给出的简便上界，不主张为最优常数。
\begin{proposition}[特定假设下的连续残差控制]\label{prop:energy}
在上述精确边界和值跳跃条件下，令$\beta_{\min}=\min(\beta_1,\beta_2)$，则
\begin{equation}
\norm{\nabla e}_{L^2(\Omega)}
\le\frac{\norm r_{L^2(\Omega)}+2^{-1/2}\norm s_{L^2(\Gamma)}}{\beta_{\min}}.
\end{equation}
其中$\norm r^2=\sum_i\norm{r_i}_{L^2(\Omega_i)}^2$。
\end{proposition}
\begin{proof}
对误差方程分部积分，得到
\[
\sum_i\int_{\Omega_i}\beta_i|\nabla e_i|^2
=\sum_i\int_{\Omega_i}r_ie_i-\int_\Gamma se.
\]
左端不小于$\beta_{\min}\norm{\nabla e}^2$，右端绝对值不大于$(\norm r+2^{-1/2}\norm s)\norm{\nabla e}$。若梯度范数为零则结论显然，否则约去该范数。
\end{proof}
这只是能量法在受限情形下的自含推导，不是本文自适应配置点算法的一般后验估计。实际网络仅在有限点弱满足外边界和值跳跃，因而不能直接把表中的离散残差代入本命题。若还存在值跳跃误差，需要构造提升并控制其导数；单独知道界面上几个点的值误差不足以完成这一步。

例如，值跳跃误差为光滑函数$a(y)$且$a(0)=a(1)=0$时，可取左侧提升为零、右侧提升$\ell_2=2(1-x)a(y)$。它在界面产生值跳跃$a$并保持外边界为零，但同时产生通量跳跃$-2\beta_2a$以及含$a''$的内部项。由此可见，处理值误差还涉及沿界面的正则性，不能只把一个值RMS与通量RMS相加就称为完整能量估计。本文联合指标是一种可计算的启发式量，其数学作用应限定在这个层面。

\section{随机特征空间、可辨识性与观测范数}
\subsection{随机与线性的两个含义}
固定隐藏参数之后，函数集合
\[
V_N=\operatorname{span}\{\phi_1,\ldots,\phi_N\}
\]
是有限维线性空间。随机性来自生成这个空间的参数，而不是输出系数求解本身。对于给定空间和给定点集，线性PDE对应确定的矩阵问题；不同随机种子改变空间以及点位置，实验离散程度因此包含不止一种随机来源。共享特征能消除方法间的一部分差异，但不能让不同种子变成完全相同的近似问题。

特征个数不一定等于空间维数。若两个特征完全相同，其系数可以相互抵消；即使数学上线性独立，也可能在有限精度或特定配置点上近乎相关。有效秩衡量的是已加权观测矩阵可以分辨的方向，它既不同于形式上的特征数，也不同于连续函数空间的精确维数。这解释了为什么本文同时报告总列数和截断后的数据秩。

采用分片空间时，$V_N=V_{N_1}^{(1)}\oplus V_{N_2}^{(2)}$中的函数在界面可以不连续。界面行把两个空间连接起来，而不是预先要求所有基函数连续。对非零值跳跃，这种表示非常直接；对零跳跃，连续性由方程约束实现。若删除界面行，两个子域求解将失去所需的耦合，内部残差小也不能单独恢复原模型。

\subsection{有限维观测控制的必要条件}
设$L$把$V_N$中的函数映射为所有加权内部、边界和界面观测。只有$L$在所考虑空间上为单射时，$\norm{Lv}$才是一个范数；若存在非零$v$使所有观测为零，它只能是半范数。有限点配置的稳定性因此首先涉及可辨识性，而不是简单地把点数加到大于特征数。

在系数表示下，若$B$满列秩，对任何预测矩阵$P$有
\begin{equation}\label{eq:prediction_bound}
\norm{Pc}_2\le\norm{PB^\dagger}_2\norm{Bc}_2
\le\frac{\norm P_2}{\sigma_{\min}(B)}\norm{Bc}_2.
\end{equation}
第一步使用$c=B^\dagger Bc$。该式是固定空间、固定观测和固定验证点的确定性结论；它提示预测稳定性常数$\norm{PB^\dagger}$比仅看$\kappa(B)$更接近具体函数预测任务，但本文尚未对全部实验计算这个常数。

若矩阵截断后，只能对保留右奇异空间上的系数使用对应界。丢弃空间中的成分不受这个训练观测界的直接控制，第二章的误差分解已将其单列。相对阈值改变时，保留空间可能变化，不能假设所有阈值下比较的是同一个近似对象。这也解释了为什么截断与配置点的交互必须结合奇异值实际检查。

\subsection{有限点小残差不等于连续小残差}
一个初等反例能说明有限观测的逻辑边界。给定区间内有限点$x_1,\ldots,x_m$，多项式$\psi(x)=\prod_{j=1}^m(x-x_j)$在这些点全部为零，但在点集之外通常不为零；任意倍数$C\psi$仍具有零经验观测，连续范数却随$|C|$增大。这个反例不是说本文tanh空间一定包含该多项式，而是说明若缺少空间限制及稳定观测假设，不能从有限点数值推出一般连续结论。

对于固定有限维特征空间，充分覆盖和矩阵可辨识性可以排除某些这种自由度，但常数仍可能很大。对于不断改变的特征或自适应点集，则需进一步控制空间和采样的共同变化。本文没有证明随机点集以何种概率满足统一范数等价，因此所有采样结论以独立离散验证和重复种子为证据，不写成未建立的概率收敛定理。


\chapter{局部随机神经网络离散方法}
\section{随机特征与物理导数}
每侧用$N_i$个固定特征$\phi_{ij}$，令$u_i^N=\sum_{j=1}^{N_i}c_{ij}\phi_{ij}$。隐藏权重与偏置随机生成后不优化。直线子域先映射到单位参考方形，左/右侧$\xi=(x-x_{\min})/.5,\eta=y$，特征为
\begin{equation}
\phi_j=\tanh\{w_{j1}(\xi-.5)+w_{j2}(\eta-.5)+b_j\}.
\end{equation}
若$t=\tanh z$，则$\sigma'=1-t^2$、$\sigma''=-2t(1-t^2)$。物理梯度带$x$方向因子2，物理Laplacian为$\sigma''(4w_{j1}^2+w_{j2}^2)$。忽略局部坐标尺度会产生模型错误，故程序以独立中心差分检查链式导数。

隐藏参数取指定均匀范围，不假设不同宽度的特征严格嵌套。相同种子改变$N_i$时，权重和偏置的随机调用顺序可能不同；宽度实验因此是同分布统计比较，不构成嵌套空间的单调性证明。

\section{分块矩阵组装}
令$c=(c_1^T,c_2^T)^T$。内部行分别为$(-\beta_1\Delta\Phi_1,0)$与$(0,-\beta_2\Delta\Phi_2)$，外边界仅激活所属侧的特征块。界面两类行是
\begin{equation}
A_u=(-\Phi_1,\Phi_2),\qquad
A_q=(-\beta_1\partial_n\Phi_1,\beta_2\partial_n\Phi_2).
\end{equation}
通量必须含各侧$\beta$，不能用导数连续代替通量条件。界面端点或角点可能造成多组约束重合，初期采用段中点或开区间随机点以避免重复。

各组样本数为$m_k$、平方权重为$\lambda_k$、固定物理尺度为$s_k$时，目标写为
\begin{equation}\label{eq:weighted}
\min_c\sum_k\frac{\lambda_k}{m_ks_k^2}\norm{A_kc-d_k}_2^2.
\end{equation}
叠置时矩阵和数据同时乘$\sqrt{\lambda_k/m_k}/s_k$。前期直线模型取边界与界面平方权重50；raw尺度全1，balanced内部尺度$\beta_i$、通量尺度$\beta_{\max}$。这些选择改变经验目标中的相对重要性，不是无损等价变换。高对比度时归一化通量小仍可对应较大物理通量误差，所以两种指标都需报告。

\subsection{分组目标与权重的含义}
设内部残差向量为$r_1,r_2$，每侧边界向量为$r_{B,i}$，两类界面向量为$r_u,r_q$。直线问题的目标可具体写为
\begin{align}
J(c)={}&\frac{\norm{r_1}^2}{m_1s_1^2}+\frac{\norm{r_2}^2}{m_2s_2^2}
+\sum_{i=1}^2\frac{\lambda_B\norm{r_{B,i}}^2}{m_{B,i}}\nonumber\\
&+\frac{\lambda_\Gamma\norm{r_u}^2}{m_\Gamma}
+\frac{\lambda_\Gamma\norm{r_q}^2}{m_\Gamma s_q^2}.
\end{align}
代码把每侧边界作为同点数的独立组归一化，两组的合计贡献在不同策略间保持一致。$m_\Gamma$同时进入两类界面行，不能把界面点数增加误解为界面平方权重也按同一比例增加。

行乘数$\gamma$进入平方目标后是$\gamma^2$，因此花形近似复现的行乘数50与直线模型的平方权重50不是同一个设置。不同残差可能对应不同物理单位；本文几何坐标无量纲，尺度选择是明确的数值约定，尚未建立跨物理模型的普适量纲平衡原则。正文同时报告物理通量，是为了防止归一化数值遮蔽实际界面偏差。

\subsection{约束数量与近似空间}
通常需要足够行数支撑$K$个系数，但$m_{\rm eq}\ge K$只是满列秩的必要条件，不是充分条件。随机特征近相关、某侧覆盖不足或重复约束都可能使矩阵近秩亏。另一方面，增加点数没有改变固定特征空间本身，不能创造网络缺少的局部尺度。新增约束如何改变解及条件数还取决于误差分布、权重和截断，见命题\ref{prop:rows}。

\section{求解与独立评估}
采用SVD型最小二乘避免显式正规方程，报告未截断条件数、保留条件数、有效秩、最小保留奇异值及系数范数。非零岭参数通过增广$[B;\sqrt\alpha I]$求解，数据矩阵秩与增广求解秩分开记录。

程序正常且指标有限仅说明运行完成，不代表精度验收或理论收敛。解误差在独立点上计算；均匀等权点的相对离散误差为$\sqrt{\sum|u_N-u|^2/\sum|u|^2}$，并报告最大误差。PDE残差使用解析算子在独立点计算。训练损失可能很小而边界、界面或未采区域误差仍大，有限点零残差不能保证连续PDE被满足。

命题\ref{prop:finite_error}给出一个解释框架：小训练目标只反映已选行及其权重，并不自动控制验证近似项、截断丢弃项或预测映射。本文测试点与训练点分离、另报物理通量及秩，正是为了降低只按一个训练数字判断算法的风险。验证点也只构成离散观测，不能代替完整连续误差分析。

\section{有限差分导数模式与原文区别}
对每个子域的平滑特征采用
\begin{align}
\partial_x\phi&\approx\frac{\phi(x+h_1,y)-\phi(x-h_1,y)}{2h_1},\\
\partial_{xx}\phi&\approx\frac{\phi(x+h_2,y)-2\phi(x,y)+\phi(x-h_2,y)}{h_2^2}.
\end{align}
即使点靠近界面，也只对同一网络特征延拓求差分，不切换到另一侧网络。对足够光滑函数，两种中心格式的截断误差为$O(h^2)$；二阶差分中浮点相减误差会除以$h_2^2$放大，因此步长并非越小越好。

花形近似复现按原文使用物理坐标$z=w\cdot x+b$及行权重$\gamma=50$\cite{interfacepreprint}，没有沿用直线模型的参考坐标、样本数归一化或平方权重50。其初始化及积分细节见第五章。本文强形式实现不是DG或PG方法，也不将解析导数版本称为原文FD严格复现。

\section{从数据到矩阵的实现细节}
\subsection{局部坐标尺度与链式法则}
设子域矩形中心为$x_i^c$，两个方向边长为$h_{i,x},h_{i,y}$，局部坐标为$\xi=(x-x_i^c)/h_{i,x}$和$\eta=(y-y_i^c)/h_{i,y}$。对$z=w_x\xi+w_y\eta+b$，特征的物理导数为
\[
\partial_x\tanh z=\frac{w_x}{h_{i,x}}(1-\tanh^2z),\quad
\partial_{xx}\tanh z=-2\frac{w_x^2}{h_{i,x}^2}\tanh z(1-\tanh^2z).
\]
另一方向同理。这里的中心化坐标与前文单位方形坐标减去$1/2$等价；直线子域取$h_{i,x}=.5,h_{i,y}=1$。若遗漏边长倒数，内部算子和通量均会改变；即使一个对称算例偶然给出较小误差，算法也不再求解原方程。本文把链式因子放在特征模块统一处理，界面组装不再重复缩放。

局部坐标映射与权重范围共同确定物理变化尺度。把同一个随机参数范围用于大小不同的子域，会产生不同的物理频率；因此报告scale时还必须说明是否映射到参考矩形。花形近似复现采用原物理坐标，不能直接把其参数1.6与直线局部映射中的scale2视为相近设置。两者即使都用tanh，矩阵列和导数的幅值也可能不同。

激活函数接近饱和时，一阶和二阶导数趋小，可能降低某些PDE观测的分辨能力。权重过小则多个特征在有限区域上接近低阶展开，容易近相关。这些都是范围敏感性的机制解释，本文没有据此推导某个普适最优范围；具体选择仍由事先登记的模型和扫描结果说明。

\subsection{装配维度、数据顺序及符号核对}
输出系数按两个子域拼接为$c=(c_1^T,c_2^T)^T$。左内部行只作用于$c_1$，右内部行只作用于$c_2$；边界按所在子域赋值。界面值行是$(-\Phi_1,\Phi_2)$，通量行为$(-\beta_1\partial_n\Phi_1,\beta_2\partial_n\Phi_2)$。右端分别放置$g_u$和$g_q$。若更换跳跃定义，整行和右端同时变号不改变约束；只改其中一侧则会求解另一个问题。

内部行、边界行及界面行分别保留组信息后再叠置，这允许单独计算每组残差并核对行数。所有组必须使用相同的系数列顺序；切分预测系数时也使用相同索引。仅检查最终矩阵可乘并不足以发现子域列交换，因为这种错误往往仍有合法维度。

制造解检验包含两层。第一层在任意非配置点计算解析PDE和跳跃，检查模型数据自身是否一致。第二层将特征导数与独立差分近似比较，检查链式因子和符号。模型一致性正确不意味着特征模块正确，反之亦然，因此两类测试分开。通量不恒零及非零跳跃的构造，可使遗漏扩散系数或漏写一侧约束更容易暴露。

\section{QR、SVD与正则化的不同作用}
\subsection{满列秩时的QR解释}
若$B$满列秩，薄QR分解为$B=QR$，其中$Q^TQ=I$且$R$可逆。将数据分解为$QQ^Td+(I-QQ^T)d$，两部分正交，得到
\[
\norm{Bc-d}^2=\norm{Rc-Q^Td}^2+\norm{(I-QQ^T)d}^2.
\]
第二项与$c$无关，因此最优系数由上三角方程$Rc=Q^Td$给出。这一推导说明QR避免形成$B^TB$；但当$R$本身接近奇异时，QR也不能凭空消除数据敏感性。实际秩亏问题可使用带列选主元的QR或SVD，本文采用SVD型实现以便统一报告奇异值和截断。

SVD把相关方向明确分解，允许以阈值舍去难以可靠区分的成分。然而阈值舍弃的可能是噪声放大方向，也可能包含逼近真解所需的信息。因此数值秩是求解约定与矩阵谱共同确定的结果，不能只按一次阈值所得秩给出特征空间的绝对数学维数。

\subsection{岭的偏差与残差权衡}
对$\alpha>0$，岭解满足
\[
c_\alpha=\sum_j\frac{\sigma_j}{\sigma_j^2+\alpha}(q_j^Td)v_j.
\]
与不截断伪逆相比，第$j$个方向乘以$\sigma_j^2/(\sigma_j^2+\alpha)$，小奇异方向收缩更强。若数据恰为$Bc_*$，则在可辨识方向上仍存在收缩偏差；如果$c_*$在这些方向上具有重要分量，训练稳定性改善也可能伴随更大的真解误差。实验中强岭导致误差增加，与这种偏差机制相容，但并不能反推出全部误差都由岭引起。

代码通过增广矩阵$[B;\sqrt\alpha I]$求解岭目标。增广矩阵的秩与原数据矩阵的有效秩不同；前者因正则化可变为满列秩，后者仍说明原观测可以分辨多少方向。论文所列data rank针对原数据矩阵，避免把正则化增加的对角行误认为采样信息增加。

\subsection{有限精度与实际求解的边界}
双精度机器精度约为$2.22\times10^{-16}$。当矩阵条件数接近其倒数时，最坏方向的相对扰动可能非常敏感；但这个判断是风险尺度，不能直接预测每个右端的实际误差。数据是否激发最坏方向、截断是否舍弃该方向，以及预测矩阵是否放大对应系数，都会影响最终函数变化。

评价算法时需要区分组装误差、线性求解误差和函数求值误差。解析特征导数减少了差分截断，但仍有浮点运算；SVD型最小二乘避免正规方程，不保证高度相关列下的所有系数准确；大系数相消又可能使函数值的实际可辨位数减少。本文固定矩阵右端实验只研究其中一个环节，未完成对全部误差来源的统一量化。

\section{有限差分特征导数的误差机制}
对足够光滑的一维函数，中心二阶差分的截断误差为$h^2\phi^{(4)}(\xi)/12$。在简单浮点误差模型中，若函数值含有量级$\epsilon_{\rm mach}|\phi|$的扰动，除以$h^2$后会出现量级$\epsilon_{\rm mach}|\phi|/h^2$的放大。一个常见的解释性模型是
\[
E(h)\lesssim C_1h^2+C_2\epsilon_{\rm mach}h^{-2}.
\]
这里常数依赖特征、导数范围和求值过程，符号$\lesssim$表示机制尺度，并非本项目全部特征上已经认证的统一上界。两项平衡提示最优尺度依赖浮点精度及函数高阶导数，而非步长越小越好。

对tanh特征，高阶导数还包含权重的相应幂，因此不同参数范围可改变差分误差。进一步把差分误差放入最小二乘矩阵后，其对解的作用又受奇异值和截断影响。仅比较单个导数近似误差不足以完全解释PDE解误差；本文同时保留解析导数对照和步长扫描，正是为了避免只观察某一环节。

差分邻点可以跨越几何界面，但对局部网络进行差分时，求值对象始终是同一侧定义的光滑特征延拓，不能在邻点位置改用另一侧网络。否则差分混合了两侧物理不连续性，不再是该侧特征导数近似。花形实现使用各自网络的物理坐标特征，内部组只在对应区域布点，界面仍使用独立两侧迹值。


\chapter{配置点策略与预算分配}
\section{基线及公平比较}
固定隐藏特征后，配置点决定哪些残差和约束被经验目标看到。本文比较A固定完整矩形网格、B固定随机点、C内部PDE残差增量点及D区域与界面联合指标。A/B一次求解；C/D先以共同初始点求解，再在同一候选池中新增点并重求输出系数，隐藏参数保持不变。

在同一随机种子内，特征流、点组流、候选流和抽样流分别管理。B的最终随机点包含C/D初始点前缀；A使用规则网格，不宣称与随机初始集相同。所有方法使用同一外边界点、固定权重和最终测试点。五种子为11、23、37、51、71，所有运行及失败状态均保存，不能只取最好一次。

界面一个物理点提供值与通量两行，故
\begin{equation}
m_{\rm eq}=m_{I,1}+m_{I,2}+2m_\Gamma+m_B.
\end{equation}
固定物理点预算时D可改变方程数；固定方程预算时D可改变物理点数。本文分别试验并同时报告两种计数，避免把新增界面方程的收益误归因于选点位置。

\section{残差指标和候选抽样}
给定当前解，内部残差$r_i=-\beta_i\Delta u_i^N-f_i$，界面残差$r_u=[u^N]-g_u$、$r_q=[\beta\partial_n u^N]-g_q$。在固定尺度下定义
\begin{equation}\label{eq:indicators}
\eta_i=\operatorname{RMS}(r_i/\beta_i),\qquad
\eta_\Gamma=\operatorname{RMS}\sqrt{r_u^2+(r_q/\beta_{\max})^2}.
\end{equation}
这些是采样指标而非已证明的后验解误差估计器。高对比度归一化可能弱化某些物理量的重要性，故实验仍报告未归一化通量。

C按残差分布选择内部新增点，但组预算和界面新增数固定；D根据三个指标调整预算，再按各组位置残差选点。可用候选点上的概率为
\begin{equation}
p_\ell=\frac{.2}{M}+.8\frac{|r_\ell|}{\sum_j|r_j|},
\end{equation}
界面用式\eqref{eq:indicators}中相应点指标代替$|r_\ell|$。全零时退回均匀概率，抽样无放回。这种探索混合避免只集中在少数候选点，但不保证连续区域覆盖或最佳配置。

这里$p_\ell$是抽取一个新增批次前，在当前可用候选集合上指定的相对权重。无放回选择后的点相互依赖，后续选择还要在剩余点间条件化；不能把批次内每点都视为按相同$p_\ell$独立抽取。最低配额保证每组获得一些点，均匀探索保证每个可用候选有正的初始选择权重，两者作用不同，也不构成最优预算定理。

\subsection{经验目标与积分测度的区别}
非均匀点改变经验目标所偏重的位置。为了说明其中区别，先考虑一个与实际增量算法不同的简化情形：固定残差函数$r$及密度$\rho$，独立有放回地抽取$X_1,\ldots,X_M$，且$\int r^2\rho<\infty$，则
\begin{equation}\label{eq:iid_sampling_expectation}
\mathbb E\left[\frac1M\sum_{j=1}^M r(X_j)^2\right]=\int r(x)^2\rho(x)\,dx.
\end{equation}
这是由期望线性性和同分布直接得到的恒等式。如果目标是另一测度$\rho_0(x)\,dx$，在$\rho>0$覆盖$\rho_0$支集且相应量可积时，单次带权变量$r(X)^2\rho_0(X)/\rho(X)$的期望才为$\int r^2\rho_0$。

本文既没有这种逆密度权重，也不是式\eqref{eq:iid_sampling_expectation}中的固定密度独立抽样：每轮残差依赖上一轮解，使用无放回候选选择，并保留历史训练点。故不能把最终目标严格等同于一个固定$\rho$下的连续积分，或直接套用iid统计保证。能明确陈述的是，实际求解始终最小化已选点上的分组经验均方，采样策略改变了这些点的位置与数量。残差概率及增量/替换采样思想已有先行研究\cite{rar}；本文研究简单联合预算机制的数值表现。

\section{最低配额及整数预算}
每轮最低新增点$(16,16,8)$，成本向量在点预算下为$(1,1,1)$，在方程预算下为$(1,1,2)$。扣除最低配额后，剩余成本按$\eta_k+\varepsilon$比例分配，除以各组点成本得到实数配额，先取整，再以剩余成本和最大欠配额补齐。成本1组保证此实验中预算能用尽。

初始内部点128/侧、界面32。静态增量为$(80,80,32)$，对应每轮192点或224标量行。主平滑比较三轮，A/B/C最终内部368/侧、界面128；峰值比较增加六轮设置，并固定最终总行1232或1904，外边界240行不变。

完整流程为：\begin{enumerate}
\item 固定模型、特征、边界、预算、候选池及随机流，保存配置。
\item 使用初始点求解；在候选池计算当前残差，不使用真解误差选点。
\item 计算组指标及整数配额，按探索混合概率新增未使用点。
\item 重新组装并求输出系数，重复预设轮数，不按最终测试误差挑轮次。
\item 在从未参与选点的网格评估全部指标，保留配额轨迹、系数及完整时间。
\end{enumerate}

\section{消融、候选偏差及复杂度}
去PDE或去界面指标时，对应组指标不参与动态预算，位置抽样退回均匀，最低配额和总预算仍不变。这个消融会同时改变配额与位置，因此后续固定比例扫描是必要的混杂因素对照，不能把消融变好直接解释为某指标无价值。

原D下一轮组RMS仍包含已经选入训练的候选点，可能乐观。D\_unused仅排除这些点后计算组RMS；位置选择仍无放回。该未用池也受历史选择影响，不是独立同分布最终测试集。本文没有据此使用iid采样的理论保证。

\subsection{三种效应的区分}
评价D时需要分开三个问题：组配额是否合理、给定配额后的残差选点是否有益、固定特征是否有能力表示需要解析的结构。主比较混合了前两个效应；去指标消融还改变位置分布。因此本文增加同最终配额的固定随机对照，并通过特征宽度和独立峰值暴露第三个问题。若在同配额下固定随机已经更好，不能把动态配额消融的收益归因于残差位置选择。

固定比例对照仍不是严格的单因素因果试验：增加界面点也减少了内部点，候选池抽样与直接随机生成的具体位置不相同。该对照用于检验“残差选点是获得该收益的必要条件”这一强解释，而不是证明界面点数为唯一原因。其参数扫描在主实验后登记，结果按探索性证据处理。

总特征$K$、矩阵行$m$、候选点$M$时，固定维度下组装约$O(mK)$，稠密SVD在$m\ge K$时主要为$O(mK^2)$，候选评估约$O(MK)$。$T$轮增量方案包含$T+1$次求解及$T$次候选评估，必须与静态的一次求解比较完整成本。重复奇异值诊断也纳入本项目计时；不能只展示最终一次求解的时间。

\section{采样分配的完整解释}
\subsection{组指标与位置指标分别起什么作用}
联合策略存在两级决定。组指标先决定下一轮新增资源在左内部、右内部和界面之间如何分配；位置指标再决定每组内部选哪些候选点。组平均残差较大，并不意味着该组内所有位置都同样困难；某个位置残差很大，也不意味着把整轮预算全部给该组最合适。这两级结构需要分开记录，否则最终误差无法对应到具体采样行为。

本文组指标使用RMS，既考虑高残差幅值，也受高残差覆盖比例影响。少数极端点可能提高指标，而一个范围较广但中等残差的区域也可能具有较大RMS。将其改为最大值、分位数或平方均值都会改变分配规则。本文未把这些备选逐一优化，只固定当前定义进行检验；不把实验后想到的指标写成已经测试的方法。

界面指标把值误差与按最大系数归一化的通量误差组合。它具有简单、无额外训练的优点，但隐含一个数值优先级：对于高系数侧，较大的物理通量误差可能被缩小到与值误差相近的量级。实验因此同时报告归一化和物理通量，不能用前者很小替代后者满足应用精度要求。

\subsection{预算分配的代数形式}
设三组指标为$\eta_k$，最低点数为$b_k$，每个物理位置成本为$a_k$，本轮成本预算为$B_0$。先令剩余成本$R=B_0-\sum_ka_kb_k$。当$R\ge0$且指标不全近零时，实数新增目标为
\[
t_k=\frac{R(\eta_k+\varepsilon)}{a_k\sum_j(\eta_j+\varepsilon)}.
\]
先取整数$\lfloor t_k\rfloor$，再在剩余成本可容纳的组间补齐。若指标均接近零，改用相等组权重，避免除以微小总量。实际代码中的阈值为$10^{-14}$，属于数值保护而非误差容限。

成本为1的组使本实验整数预算总能填满；对任意成本向量并不具有这个保证。例如所有可选成本均为2而剩余成本为1时无法填满，必须改成不超过预算或重新定义允许成本。本文算法对本研究两种成本约定合法，不宣称解决一般整数规划问题。

最低配额会在指标非常不平衡时限制集中程度，避免某侧长期完全不新增点，但也会消耗可用于困难区域的预算。这一选择是覆盖与集中之间的工程折中。既有数据没有证明$(16,16,8)$是最优配额，也没有考察把它改成随轮次变化的规则；应按真实设置理解结果。

\subsection{抽样概率并不等于最终点密度}
探索比例0.2使可用池内每点在当前批次开始时具有非零选择权重。它不表示最终训练集中恰有20\%的点均匀分布，因为无放回过程、历史点、组配额和多轮更新共同决定最终集合。同样，最终点数多的区域不一定在每轮都具有最大的残差。

历史点保留可避免丢弃已经满足的局部约束，但也使早期选点持续影响后续目标。若采用替换策略，每轮更新后的训练测度和求解成本会不同。本文只实现增量路线，并在比较时把所有历史求解计入时间。不能把最终一次拟合时间当成完整自适应成本，也不能把与增量路线不同的已有采样方法视为逐公式复现。

候选池大小还限制能发现的结构。若窄峰附近没有足够候选点，残差排序无法选择池外位置；若池内存在相关或近重复位置，抽样无放回也不等于几何上有良好间距。均匀探索降低极端集中的风险，但不提供填充距离或连续覆盖保证。局部峰失败因此需要同时考虑候选分辨率与特征空间，不能只讨论界面预算。

\section{实验设计与结论强度}
\subsection{训练、候选、监测和最终测试}
训练点决定矩阵，候选点决定自适应更新，最终测试点只计算误差。三者用途不同。候选池虽在最初生成时与初始训练分开，其中部分位置后来会被选入训练；随着算法推进，候选指标已经受到选择历史影响，不能继续把它当作未接触训练的泛化测试。

D\_unused排除已经选用的候选位置后计算组指标，有助于诊断训练点混入带来的乐观性，但未用池仍经过历史选择筛选，故也不是新的独立样本。一个更严格的后续方案是在每轮另生成仅用于监测的池，且不从中挑点；若又根据监测误差调参数，则还需要保留未用于选择的最终测试。本文没有把这一后续设计冒充已经执行的实验。

现有测试网格与采样点的分离避免了最直接的训练误差自评，但探索性研究反复查看同一组测试结果仍可能形成研究者层面的选择偏差。因此本文明确区分主比较与后续诊断，不把后续扫描中表现较好的参数当成主实验事先确定的最优方案。真正的外推检验需要新模型或冻结规则后的新增运行。

\subsection{共享随机种子与统计单位}
同一设置内，A、B、C、D使用相同种子生成特征，这使配对差异较少受到完全不同空间的影响。不过每种策略的点集仍不同，种子共享只构成控制条件之一。跨系数组合重复相同种子时，相应特征空间高度相关，不能把三组五种子简单当作十五个独立空间。

本文以每个固定模型、预算和参数组合为一个设置，在其中报告五次重复及配对比值。跨设置合并的中位数只作为描述摘要，不附上假设独立重复后才成立的显著性结论。五个种子也不足以准确估计罕见失败概率；没有看到失败不等于已证明失败概率为零。

标准差和均值的用途不同。均值反映该五种子集合的平均误差，标准差反映离散程度，最大值揭示是否存在明显较差初始化。对于数值跨度较大的误差，少数较差运行会影响算术均值；本文保留逐次CSV，并同时使用配对比例中位数，避免只呈现一种摘要而失去分布信息。

\subsection{结果判据不随预期移动}
正确性判据包括制造数据一致、导数一致、FDM收敛与预算守恒。自适应效果判据则是合理基线、公平资源、独立误差和完整成本，不能要求策略必胜，也不能在其失败后悄悄换指标宣称成功。PDE残差稍低但真解误差更高时，应报告为取舍，不把原先的解精度目标事后改成仅优化残差。

数值状态ok只表示程序运行完成并返回有限指标。窄峰误差超过1说明当前配置没有达到高精度，即使最小二乘本身成功退出。本文把运行异常、低精度与未达到改进目标分别表述，因为三者的修正措施不同：异常需要程序排查，低精度需要模型或离散诊断，缺乏优势则可能说明当前方案本身不值得继续加复杂度。

\section{计算成本的可比较范围}
相同物理点数只控制一种资源，界面位置产生两行的事实使行数仍可能不同；相同行数又不控制几何覆盖。本文同时记录二者，并指出增量算法还有候选评估和反复求解开销。因此固定最终行预算的比较回答的是相近最终线性系统规模下的效果，不等价于严格相同总CPU时间。

在CPU稠密线性代数实现中，矩阵宽度增长通常比少量新增行更昂贵，奇异值诊断本身也有成本。本文计时包含这些诊断，是可复核的实现约定；若将来只用QR而不重复诊断，运行时间可能改变，但必须对全部方法一致调整。当前不能仅凭复杂度阶或某次耗时推断另一台机器的加速倍数。

进一步的精度目标比较应预先规定误差阈值、允许参数和搜索成本，报告达到阈值所需的完整时间。现有实验主要为固定配置与预算比较，尚未完成统一精度目标下的完整优化搜索。本文因此没有声称已经证明某方法在任意目标精度下更高效，也没有用不同算法各自最好的一次运行代替公平成本分析。


\chapter{数值实验与结果分析}
\section{环境、指标与数据来源}
实际运行环境为Linux CPU、Python 3.12.14、NumPy 2.2.6、SciPy 1.15.3及Matplotlib 3.10.3，BLAS线程设为1。程序设计为Windows Python 3.12可运行，但尚未完成Windows实机验证。无需GPU或真实外部数据。直线模型多数使用五种子，花形近似复现主比较使用十种子；表中的均值覆盖全部指定运行，样本标准差采用分母$n-1$。

\begin{table}[htbp]\centering\small
\caption{主要实验设置概览，细项以各run配置为准}\label{tab:protocol_overview}
\begin{tabular}{lp{.69\textwidth}}\toprule
设置 & 固定参数与比较范围\\\midrule
平滑采样主比较 & 每侧64特征，scale2；三轮；A/B/C/D及两消融；三系数组合；点预算1104或行预算1232；五种子。\\
局部峰比较 & 每侧128特征，scale2/4；峰宽.12/.06；三轮/六轮；1232/1904行；两个系数组合；五种子。\\
花形主近似复现 & 每侧80/160/320特征；总配置点1000/1000/5000；两网络权重范围1.6、偏置.7；FD和解析导数；十种子。\\
\bottomrule\end{tabular}\end{table}

直线模型测试每侧83乘83中点及257界面点；峰值细测试每侧161乘161、界面521点，另用81乘81检查积分敏感性。全局Poisson随机特征使用81乘81内部点。FDM误差在其节点上，残差为离散系统相对残差；其他PDE残差为独立点上的连续解析算子残差，两者不能直接作为同指标比较。时间区分单次组装/求解和完整自适应过程，不含最终验证/绘图。

所有表格由原始CSV生成，图由CSV重绘。主要来源为实验01、02、03、04、06、07、08及09的两组复现记录，均为2026年10月13日保存。完整路径见附录；原文数据另标注，不混为本项目输出。

设同一设置下五次误差为$e_s$，报告均值与样本标准差：
\[\bar e=\frac1n\sum_s e_s,\qquad s_e=\sqrt{\frac{\sum_s(e_s-\bar e)^2}{n-1}}.\]
标准差描述随机初始化的离散程度，不等于均值的置信区间。采样比较采用共享种子的配对比例$R_s=e_{D,s}/e_{X,s}$，其中$X$为基线；$R_s<1$仅表示该次解误差更低。比例中位数与均值之比不同，本文不混用。五种子不足以支持广泛的统计显著性或总体优势声明。

\section{Poisson基线正确性}
表\ref{tab:fdm}和图\ref{fig:fdm}显示四级网格的误差与观测阶。相邻步长减半后误差约降四倍，与第二章的正弦特征值推导一致。单元测试还核查矩阵符号、制造右端、tanh导数和秩亏最小范数解。
\begin{table}[htbp]\centering\caption{Poisson五点差分真实结果}\label{tab:fdm}\begin{tabular}{lll}\toprule
每向内部点 & 相对离散误差 & 观测阶 \\ \midrule
15 & $3.219\times10^{-3}$ & -- \\
31 & $8.036\times10^{-4}$ & 2.0021 \\
63 & $2.008\times10^{-4}$ & 2.0005 \\
127 & $5.020\times10^{-5}$ & 2.0001 \\
\bottomrule\end{tabular}
\end{table}
\figurefile{fdm_convergence}{正弦Poisson五点差分收敛；数据来自exp01原始CSV。}{fig:fdm}
基础全局RaNN每种特征数五次：64特征相对误差均值约$8.64\times10^{-3}$，128特征约$5.99\times10^{-5}$，后者样本标准差约$3.08\times10^{-5}$。与此同时128特征加权矩阵平均条件数约$2.48\times10^9$。精度改善并不意味着条件性改善。全局网络与双子域网络总特征数、坐标尺度、配置点及评估方式不同，本文不据此宣称公平效率优越。

\section{直线界面与尺度取舍}
150次运行包含连续与非零跳跃、七组正反向系数组合、raw/balanced及Poisson极限。无运行异常，平衡组通过预设粗误差门槛；这一门槛不覆盖后续窄峰。表\ref{tab:interface}为非零跳跃、每侧128特征、scale2的五次均值。
\begin{table}[htbp]\centering\caption{高对比度下解与物理通量，五种子均值及样本标准差}\label{tab:interface}\resizebox{\textwidth}{!}{\begin{tabular}{llll}\toprule
$(\beta_1,\beta_2)$ & 尺度 & 相对$L^2$：均值$\pm s$ & 通量RMS：均值$\pm s$ \\ \midrule
(1000,1) & raw & $(4.579\pm4.578)\times10^{-5}$ & $(1.763\pm2.570)\times10^{-4}$ \\
(1000,1) & balanced & $(1.391\pm0.681)\times10^{-7}$ & $(2.756\pm3.733)\times10^{-4}$ \\
(1,1000) & raw & $(1.094\pm0.898)\times10^{-4}$ & $(1.711\pm1.554)\times10^{-4}$ \\
(1,1000) & balanced & $(1.466\pm0.280)\times10^{-7}$ & $(4.453\pm2.523)\times10^{-4}$ \\
\bottomrule\end{tabular}
}\end{table}
\figurefile{interface_contrast}{非零跳跃、每侧128特征的系数对比度实验；物理通量未除以系数。}{fig:interface}
平衡后解误差明显降低，但物理通量均值不一定降低。raw最高条件数约$1.69\times10^{15}$、最低有效秩240/256；balanced仍可达到约$2.11\times10^{12}$、最低255/256。截断SVD使近秩亏系统可求解，但不能据正常运行声称病态已消除。

\section{采样主比较与消融}
180次主比较固定每侧64特征、scale2、三轮；每种预算覆盖三个系数组合和五种子。表\ref{tab:sampling}展示点预算、$\beta=(1000,1)$的全部六种方法，最终物理总点1104，D方程数可变。
\begin{table}[htbp]\centering\caption{等物理点预算下采样与消融，误差为均值及样本标准差}\label{tab:sampling}\resizebox{\textwidth}{!}{\begin{tabular}{llll}\toprule
方法 & 相对$L^2$：均值$\pm s$ & 完整时间/秒 & 界面点均值 \\ \midrule
A & $(1.271\pm0.411)\times10^{-4}$ & 0.0133 & 128.0 \\
B & $(1.200\pm0.404)\times10^{-4}$ & 0.0137 & 128.0 \\
C & $(1.555\pm0.542)\times10^{-4}$ & 0.0763 & 128.0 \\
D & $(1.591\pm0.541)\times10^{-4}$ & 0.0767 & 68.0 \\
去PDE指标 & $(1.103\pm0.253)\times10^{-4}$ & 0.0829 & 512.0 \\
去界面指标 & $(1.561\pm0.545)\times10^{-4}$ & 0.0743 & 56.0 \\
\bottomrule\end{tabular}
}\end{table}
\figurefile{sampling}{点预算下全部种子的解误差，横线为均值；数据来自exp03。}{fig:sampling}
在每种预算的3设置乘5种子配对中，D对A/B解误差胜出均为0/15，D/B误差比例中位数约1.264和1.256。设置之间相关，不能把15个配对当作15个独立种子做显著性推断。D对C仅部分改善，仍没有稳定胜出。去PDE指标消融给界面更多点而较好，但后续相同最终组配额的固定随机在六组均值中更好且更快，未支持把收益简单归因于残差选点。

这说明比较基线必须合理：若只选界面明显欠采样的静态比例，可能高估动态分配的价值。另一方面，D默认指标偏向内部，界面配额较低，也可能影响解误差。本研究把该解释保留为数值猜想，而非因果证明。

表\ref{tab:paired_sampling}按系数与预算拆开配对结果，避免只报告合并中位数而掩盖不同设置。三个系数组合均使用相同五个种子，跨设置重复并不产生新的独立初始化样本。
\begin{table}[p]\centering\caption{平滑模型的分设置配对；比值为D误差除以基线误差}\label{tab:paired_sampling}\begin{tabular}{lllll}\toprule
预算 & $(\beta_1,\beta_2)$ & 配对 & 误差较小次数 & 比值中位数 \\ \midrule
点 & (1,1) & D/A & 0/5 & 1.2605 \\
点 & (1,1) & D/B & 0/5 & 1.2619 \\
点 & (1,1) & D/C & 0/5 & 1.0175 \\
点 & (1000,1) & D/A & 0/5 & 1.2338 \\
点 & (1000,1) & D/B & 0/5 & 1.2924 \\
点 & (1000,1) & D/C & 1/5 & 1.0308 \\
点 & (1,1000) & D/A & 0/5 & 1.2721 \\
点 & (1,1000) & D/B & 0/5 & 1.2130 \\
点 & (1,1000) & D/C & 3/5 & 0.9975 \\
行 & (1,1) & D/A & 0/5 & 1.2248 \\
行 & (1,1) & D/B & 0/5 & 1.2451 \\
行 & (1,1) & D/C & 2/5 & 1.0156 \\
行 & (1000,1) & D/A & 0/5 & 1.2339 \\
行 & (1000,1) & D/B & 0/5 & 1.3013 \\
行 & (1000,1) & D/C & 1/5 & 1.0183 \\
行 & (1,1000) & D/A & 0/5 & 1.2726 \\
行 & (1,1000) & D/B & 0/5 & 1.2561 \\
行 & (1,1000) & D/C & 2/5 & 1.0197 \\
\bottomrule\end{tabular}
\end{table}

\subsection{配额与选点位置的分离诊断}
去PDE指标消融在点预算下最终每侧内部176点、界面512点，外边界240点，共1104个物理点及1616个标量方程。后续固定随机对照采用同一最终配额、特征数及随机种子，表\ref{tab:allocation_control}显示其平均解误差在三个系数组合下均更低。该比较是看到消融结果后登记的探索性诊断，不能改写成预先设计的验证。表中成本来自不同批次，支持当前环境下的成本观察；严格串行重测的证据见峰值场景。
\begin{table}[htbp]\centering\caption{相同最终配额的随机对照与去PDE指标消融，五种子}\label{tab:allocation_control}\resizebox{\textwidth}{!}{\begin{tabular}{llll}\toprule
$(\beta_1,\beta_2)$ & 同配额方法 & 相对$L^2$：均值$\pm s$ & 完整时间/秒 \\ \midrule
(1,1) & 固定随机 & $(1.112\pm0.528)\times10^{-4}$ & 0.0168 \\
(1,1) & 去PDE指标 & $(1.156\pm0.594)\times10^{-4}$ & 0.0826 \\
(1000,1) & 固定随机 & $(1.074\pm0.299)\times10^{-4}$ & 0.0169 \\
(1000,1) & 去PDE指标 & $(1.103\pm0.253)\times10^{-4}$ & 0.0829 \\
(1,1000) & 固定随机 & $(9.311\pm5.096)\times10^{-5}$ & 0.0167 \\
(1,1000) & 去PDE指标 & $(9.922\pm5.283)\times10^{-5}$ & 0.0826 \\
\bottomrule\end{tabular}
}\end{table}
这个对照未支持残差排序选点对消融收益的必要性，但也不是完整因果分解：增量历史、位置与求解路径仍不同。因各组按点数取平均，增加界面点没有直接增加界面平方权重，而是改变界面经验测度并减少内部预算。不能将此同配额实验与1232行默认方案直接称为等行预算比较。

\section{参数、秩与误差饱和}
285次探索性后续实验包含固定比例及单因素扫描。表\ref{tab:width}为$\beta=(1,1000)$宽度对照，其他设置固定。
\begin{table}[htbp]\centering\caption{特征宽度敏感性，数据秩上限为总特征数}\label{tab:width}\begin{tabular}{llll}\toprule
每侧特征 & 相对$L^2$均值 & 原始条件数均值 & 数据秩范围 \\ \midrule
32 & $1.138\times10^{-2}$ & $1.823\times10^{5}$ & 64--64 \\
64 & $1.085\times10^{-4}$ & $1.516\times10^{8}$ & 128--128 \\
128 & $1.236\times10^{-7}$ & $9.579\times10^{11}$ & 255--256 \\
256 & $2.581\times10^{-10}$ & $1.076\times10^{17}$ & 404--417 \\
\bottomrule\end{tabular}
\end{table}
\figurefile{width}{三个系数组合下的宽度、物理通量与条件数；数据来自exp04。}{fig:width}
当前平滑模型增宽改善误差，但条件数升至约$10^{17}$，256/侧时有效秩仅404至417，不能以总特征数代表独立维度。scale从.5到8的影响显著，小范围可适合平滑函数却加重病态，大范围可能使近似质量变差。界面权重、岭和点数也改变排序：强岭增加偏差；点数216至1728的解误差不单调，物理通量却可改善。这些扫描属于看到主比较后登记的探索性研究，不把其中最好参数改写为事先最优基线。

\section{独立局部峰与多预算}
在原分片制造解两侧加$h=16A x(1-x)y(1-y)\exp(-((x-c_x)^2+(y-c_y)^2)/w^2)$，取$A=.75,(c_x,c_y)=(.46,.55)$，两峰宽$.12,.06$。值跳跃不变，但通量增加$(\beta_2-\beta_1)h_x$，须重新制造数据。400次比较采用每侧128特征、scale2/4、两个行预算及五方法五种子。两个测试网格误差相对变化最大约0.49\%，未进入选点。
\figurefile{peaks}{scale4下独立峰的两个标量行预算比较；所有预设模型与方法均保留。}{fig:peaks}
80描述性配对中，D对A、B、C误差更低分别6、36、49次；对应比值中位数约1.293、1.095、.963。D\_unused没有稳定逆转结果。部分窄峰相对误差大于1，最大约1.493，无运行异常并不代表精度验收。固定特征表达力和采样目标变化都可能影响结果，不能只增加点数就承诺解决。

主峰值前段与谱任务并发，严格成本讨论使用另行25次串行重测。表\ref{tab:peaks}为固定窄峰、$\beta=(1000,1)$、scale4、六轮及1904方程，系数已与主case校验一致。D在该设置比C精度更好，但仍不如A且更贵。
\begin{table}[htbp]\centering\caption{固定峰值场景的串行完整成本，误差均值及样本标准差}\label{tab:peaks}\resizebox{\textwidth}{!}{\begin{tabular}{lll}\toprule
方法 & 相对$L^2$：均值$\pm s$ & 串行完整时间/秒 \\ \midrule
A & $(1.948\pm0.610)\times10^{-1}$ & 0.0481 \\
B & $(3.839\pm1.098)\times10^{-1}$ & 0.0450 \\
C & $(3.524\pm1.360)\times10^{-1}$ & 0.4037 \\
D & $(2.845\pm0.859)\times10^{-1}$ & 0.4097 \\
D\_unused & $(2.742\pm0.801)\times10^{-1}$ & 0.4122 \\
\bottomrule\end{tabular}
}\end{table}

窄峰结果并未唯一识别误差来源。表达空间、选点分布、组权重及数值截断均可能限制精度；尚未计算该空间的最佳近似误差，因此不能把失败严格归因于表达力。可以确定的是：在本文已运行的固定特征和预算下，联合采样没有实现稳定的精度或成本优势。

\section{固定矩阵扰动观察}
15组矩阵谱、三阈值和两个扰动方向构成90组观察。右端相对扰动固定$10^{-8}$；最小保留左奇异方向的放大与$1/\sigma_{\min,\rm retained}$在浮点误差内一致，图\ref{fig:perturb}展示系数和函数变化的差别。
\figurefile{perturbation}{固定TSVD空间下的右端扰动；理论界仅针对固定矩阵和保留空间。}{fig:perturb}
最大系数相对变化约275.6，而预测函数相对变化最大约$1.71\times10^{-8}$。冗余特征可使系数敏感而函数变化较小，但该数值还可能接近大系数抵消的浮点求值地板，未以高精度算术复核。不能由此断言病态无害，也没有考察矩阵扰动、阈值跨秩或所有物理数据误差。

\section{花形界面FD近似复现}
选Li--Wang预印本Example4.1\cite{interfacepreprint}，$\Omega=(-1,1)^2$，中心$.02\sqrt5$，极径$r=.4+.2\sin20\theta$。图1核查花内为$\Omega_1$。两侧真解为$e^{xy}/\beta_1$及$\sin x\sin y/\beta_2$，$\beta=(1,10)$；法向由曲线切向旋转并归一化，不能只用径向单位向量。

按最贴合文字的解释，两网络权重范围均为1.6、偏置范围.7。与原文相同的设置包括单隐层tanh、3:1:1点数比例、$h_1=10^{-6},h_2=5\times10^{-4}$、行权重50、GL20及十次平均。作者种子、dtype、最小二乘阈值、确切界面测度未获实现；本项目float64/gelsd/$\tau=10^{-12}$、主角度均匀并补弧长对照。因此\textbf{只称近似复现}。
\begin{table}[htbp]\centering\small\caption{原文表1与本项目近似复现，十种子均值}\label{tab:reproduction}\resizebox{\textwidth}{!}{\begin{tabular}{lllll}\toprule
$m/N$ & 导数 & 原文表1 & GL20：均值$\pm s$ & GL100：均值$\pm s$ \\ \midrule
80/1000 & fd & $1.180\times10^{-3}$ & $(3.169\pm2.090)\times10^{-5}$ & $(3.236\pm2.147)\times10^{-5}$ \\
80/1000 & analytic & $1.180\times10^{-3}$ & $(3.169\pm2.090)\times10^{-5}$ & $(3.236\pm2.147)\times10^{-5}$ \\
160/1000 & fd & $9.200\times10^{-6}$ & $(1.059\pm0.664)\times10^{-7}$ & $(1.080\pm0.683)\times10^{-7}$ \\
160/1000 & analytic & $9.200\times10^{-6}$ & $(1.059\pm0.664)\times10^{-7}$ & $(1.080\pm0.683)\times10^{-7}$ \\
320/5000 & fd & $3.170\times10^{-8}$ & $(3.057\pm0.712)\times10^{-10}$ & $(3.076\pm0.457)\times10^{-10}$ \\
320/5000 & analytic & $3.170\times10^{-8}$ & $(2.950\pm0.721)\times10^{-10}$ & $(2.961\pm0.439)\times10^{-10}$ \\
\bottomrule\end{tabular}
}\end{table}
\figurefile{reproduction}{最贴合文字初始化下的FD与解析导数对照；原文曲线不是本项目数据。}{fig:reproduction}
表\ref{tab:reproduction}的误差与原文明显不同。更低数值不能当作算法改进，差异可能涉及初始化、dtype、截断与采样，尚无法唯一归因。最初按子域解释范围的110次保留为独立对照，核对文字后追加60次权重/偏置解释，旧数据不改写。原文通量印刷矩阵疑似漏$\beta$，本项目按物理模型保留，正式版仍待核对。

在子域范围解释变体中，五种子的FD步长结果见表\ref{tab:steps}。过大步长有截断偏差，过小步长出现舍入放大；不能把此变体当作作者结论。GL20与GL100在极小误差下可有明显相对变化，均保留而不选择较好的积分值。
\begin{table}[htbp]\centering\caption{FD步长独立诊断，$m=160,N=1000$}\label{tab:steps}\begin{tabular}{lll}\toprule
二阶差分步长 & 相对$L^2$均值 & 解析PDE残差RMS \\ \midrule
$1.000\times10^{-2}$ & $2.635\times10^{-8}$ & $4.659\times10^{-6}$ \\
$5.000\times10^{-4}$ & $5.976\times10^{-10}$ & $2.213\times10^{-7}$ \\
$1.000\times10^{-5}$ & $1.262\times10^{-7}$ & $1.768\times10^{-5}$ \\
$1.000\times10^{-6}$ & $8.669\times10^{-6}$ & $8.555\times10^{-4}$ \\
\bottomrule\end{tabular}
\end{table}

定义积分敏感度为$\Delta=|e_{20}-e_{100}|/e_{100}$。表\ref{tab:quadrature}使用最贴合原文字面初始化的FD组，逐种子计算后报告中位数与最大值。GL100是更密的数值验证规则，不是连续误差的严格真值；改变规则只用于诊断，没有进入训练或选择最好运行。误差极小时较大的相对差别也可能对应很小的绝对差别，故同时报告误差量级。
\begin{table}[htbp]\centering\caption{花形界面误差对积分规则的敏感度，十种子}\label{tab:quadrature}\resizebox{\textwidth}{!}{\begin{tabular}{llll}\toprule
$m/N$ & 积分差别中位数 & 积分差别最大值 & GL100误差均值 \\ \midrule
80/1000 & 2.22\% & 3.38\% & $3.236\times10^{-5}$ \\
160/1000 & 1.52\% & 4.83\% & $1.080\times10^{-7}$ \\
320/5000 & 7.58\% & 33.30\% & $3.076\times10^{-10}$ \\
\bottomrule\end{tabular}
}\end{table}

\section{实验结论的范围}
已有数据支持对采样、尺度、表达力和求解截断进行系统讨论，但不支持普适自适应优势。平滑、峰值及花形模型具有不同的近似难度；对一个算例表现好的范围不应直接迁移。还未做传统界面FDM/FEM或PINN同预算对照、完整多层增长、复杂物理或矩阵扰动。初稿中所有结论限于已登记模型与参数，后续应以独立模型和理论补充，而非对旧测试集无限调参。

\section{多指标证据与采样过程分析}
\subsection{相同标量行预算下的独立误差}
为避免只根据相对解误差评价算法，表\ref{tab:multimetric}补列1232行主实验的最大绝对误差、两侧归一化PDE残差、界面值RMS和物理通量RMS。全部为五种子的算术均值；原始逐次数据与标准差仍保存于同一CSV及汇总表。归一化PDE残差除以各自扩散系数，而最后一列通量未作归一化，读表时不可把这些量直接按数值大小比较。
\begin{table}[p]\centering
\caption{相同1232行预算下的独立多指标均值；每侧64特征、scale2、三轮}\label{tab:multimetric}
\resizebox{\textwidth}{!}{\begin{tabular}{lllllll}\toprule
$(\beta_1,\beta_2)$ & 方法 & $E_\infty$ & $R_1/\beta_1$ & $R_2/\beta_2$ & $J_u$ & $J_q$ \\ \midrule
(1,1) & A & $4.895\times10^{-4}$ & $2.989\times10^{-3}$ & $3.125\times10^{-3}$ & $1.124\times10^{-4}$ & $9.465\times10^{-5}$ \\
(1,1) & B & $4.449\times10^{-4}$ & $3.866\times10^{-3}$ & $4.382\times10^{-3}$ & $1.185\times10^{-4}$ & $1.179\times10^{-4}$ \\
(1,1) & C & $5.999\times10^{-4}$ & $3.227\times10^{-3}$ & $3.010\times10^{-3}$ & $1.519\times10^{-4}$ & $1.236\times10^{-4}$ \\
(1,1) & D & $6.116\times10^{-4}$ & $2.974\times10^{-3}$ & $2.988\times10^{-3}$ & $1.452\times10^{-4}$ & $1.074\times10^{-4}$ \\
(1000,1) & A & $4.537\times10^{-4}$ & $3.108\times10^{-3}$ & $2.656\times10^{-3}$ & $1.123\times10^{-4}$ & $1.029\times10^{-1}$ \\
(1000,1) & B & $3.964\times10^{-4}$ & $4.613\times10^{-3}$ & $3.906\times10^{-3}$ & $1.093\times10^{-4}$ & $1.047\times10^{-1}$ \\
(1000,1) & C & $5.385\times10^{-4}$ & $3.075\times10^{-3}$ & $2.574\times10^{-3}$ & $1.450\times10^{-4}$ & $1.376\times10^{-1}$ \\
(1000,1) & D & $5.506\times10^{-4}$ & $3.036\times10^{-3}$ & $2.502\times10^{-3}$ & $1.436\times10^{-4}$ & $1.219\times10^{-1}$ \\
(1,1000) & A & $4.082\times10^{-4}$ & $2.665\times10^{-3}$ & $2.913\times10^{-3}$ & $7.398\times10^{-5}$ & $8.728\times10^{-2}$ \\
(1,1000) & B & $3.720\times10^{-4}$ & $3.402\times10^{-3}$ & $3.868\times10^{-3}$ & $7.790\times10^{-5}$ & $9.671\times10^{-2}$ \\
(1,1000) & C & $4.923\times10^{-4}$ & $2.733\times10^{-3}$ & $2.961\times10^{-3}$ & $9.955\times10^{-5}$ & $1.135\times10^{-1}$ \\
(1,1000) & D & $5.153\times10^{-4}$ & $2.655\times10^{-3}$ & $2.869\times10^{-3}$ & $9.029\times10^{-5}$ & $1.095\times10^{-1}$ \\
\bottomrule\end{tabular}
}
\end{table}

在$\beta=(1000,1)$设置中，D的右侧归一化PDE残差均值约$2.50\times10^{-3}$，低于B的约$3.91\times10^{-3}$；但D的最大绝对误差约$5.51\times10^{-4}$，高于B的约$3.96\times10^{-4}$。这提供了一个具体反例：在独立点上内部残差降低，也不能不加条件地推出最大解误差降低。原因可能同时涉及界面、边界、截断及不同误差传播方向，不能只根据这两列唯一归因。

同一设置下，D的物理通量RMS约0.122，B约0.105，说明残差目标在内部改善时并未同时改善全部界面指标。这个量级来自每侧64特征的采样主比较，不能与每侧128特征的界面验证表混用。若不注明特征数，读者可能误以为高对比度通量结果前后矛盾；实际上两表属于不同近似能力和配置点实验。

对于单位系数，D两侧PDE残差也处于较低水平，但界面值误差和最大解误差没有同步胜过静态基线。因此本研究的负面结果不仅是某个高对比度下的偶然异常。另一方面，不能把当前三个平滑设置概括为联合采样在所有问题上无效；峰值实验已经显示D相对仅PDE采样在部分设置改善，这种差异仍需按问题和指标分别报告。

\subsection{条件数较低为何仍可能误差较高}
表\ref{tab:rankcost}显示相同行预算的条件数、有效秩、物理点及时间。这里总列数为128，所有列出的主比较有效秩均为128，原始条件数约为$10^8$量级。D在高对比度设置下的平均条件数略低于B，却仍没有更小的解误差。这说明条件数反映最坏线性敏感性，不能直接作为近似误差的单调代理。
\begin{table}[p]\centering
\caption{1232行主比较的矩阵与资源诊断，五种子；时间含完整自适应过程}\label{tab:rankcost}
\resizebox{\textwidth}{!}{\begin{tabular}{lllllll}\toprule
$(\beta_1,\beta_2)$ & 方法 & $\kappa$均值 & 有效秩 & 物理点均值 & 界面点均值 & 时间/秒 \\ \midrule
(1,1) & A & $1.812\times10^{8}$ & 128--128 & 1104.0 & 128.0 & 0.0135 \\
(1,1) & B & $1.766\times10^{8}$ & 128--128 & 1104.0 & 128.0 & 0.0137 \\
(1,1) & C & $1.618\times10^{8}$ & 128--128 & 1104.0 & 128.0 & 0.0752 \\
(1,1) & D & $1.636\times10^{8}$ & 128--128 & 1168.4 & 63.6 & 0.0755 \\
(1000,1) & A & $1.314\times10^{8}$ & 128--128 & 1104.0 & 128.0 & 0.0167 \\
(1000,1) & B & $1.311\times10^{8}$ & 128--128 & 1104.0 & 128.0 & 0.0132 \\
(1000,1) & C & $1.238\times10^{8}$ & 128--128 & 1104.0 & 128.0 & 0.0773 \\
(1000,1) & D & $1.183\times10^{8}$ & 128--128 & 1168.6 & 63.4 & 0.0837 \\
(1,1000) & A & $1.507\times10^{8}$ & 128--128 & 1104.0 & 128.0 & 0.0139 \\
(1,1000) & B & $1.516\times10^{8}$ & 128--128 & 1104.0 & 128.0 & 0.0142 \\
(1,1000) & C & $1.365\times10^{8}$ & 128--128 & 1104.0 & 128.0 & 0.0752 \\
(1,1000) & D & $1.348\times10^{8}$ & 128--128 & 1169.2 & 62.8 & 0.0759 \\
\bottomrule\end{tabular}
}
\end{table}

首先，矩阵条件数不包含真解到特征空间的距离。一个条件良好但空间表达能力不足的系统，也可能给出准确求解的错误近似。其次，相同行预算不保证相同的约束经验分布；矩阵列的可辨识性略有改善，可能伴随对界面约束覆盖的改变。再次，真解误差对应特定右端及预测位置，最坏方向未必被显著激发。

这与第二章固定矩阵误差分解一致：除了扰动放大项，还存在验证近似、截断和求值项。当前主比较未截断出明显秩亏，不能据此排除空间与经验目标的影响；后续增宽实验出现的有效秩不足则属于另一种情形。统一使用条件数解释所有实验，会混淆不同宽度和截断状态。

时间方面，自适应方案约为多次求解及候选评估之和，而A/B只求解一次。表中小规模CPU时间受系统状态影响，不宜将单次差值解释为精确硬件常数；但完整成本显著增加且解误差没有改善，足以限制当前方案的实用性主张。更严格的固定峰值串行重测已单独报告，不能拿并发阶段的时间去宣称加速。

\subsection{配额轨迹提供的机制证据}
表\ref{tab:history}列出$\beta=(1000,1)$时D的三轮平均轨迹，均值不必是整数。每次实际运行中的配额均为整数并满足预算，不能把均值中的小数理解为程序使用了小数个点。
\begin{table}[htbp]\centering
\caption{联合策略D的实际平均配额轨迹；每侧64特征、五种子}\label{tab:history}
\begin{tabular}{lllll}\toprule
预算 & 轮次 & 左内部均值 & 右内部均值 & 界面均值 \\ \midrule
点 & 0 & 128.0 & 128.0 & 32.0 \\
点 & 1 & 212.4 & 225.6 & 42.0 \\
点 & 2 & 307.0 & 310.2 & 54.8 \\
点 & 3 & 399.0 & 397.0 & 68.0 \\
行 & 0 & 128.0 & 128.0 & 32.0 \\
行 & 1 & 223.0 & 238.6 & 41.2 \\
行 & 2 & 331.2 & 332.8 & 52.0 \\
行 & 3 & 438.6 & 426.6 & 63.4 \\
\bottomrule\end{tabular}

\end{table}

点预算下最终界面平均68点，少于静态128点，而内部各侧平均接近400点；行预算下界面平均约63.4点。对应物理点均值超过1104，是因为一个界面位置占两行，减少界面后可用同样行预算放更多内部位置。这一计数关系解释了表\ref{tab:rankcost}中D的物理点更多，却没有使用额外标量行。

轨迹支持“当前指标倾向把资源给内部”这一描述，但不单独证明界面点少就是误差变大的唯一原因。改变配额同时改变位置分布、组经验平均和拟合结果；还可能存在某些内部点重复提供近似信息。后续同配额固定随机对照只回答一项有限问题，不能完成全部机制识别。

如果以后改进分配规则，首先应规定如何平衡物理通量与归一化通量，以及是否按每组单位成本的收益分配。当前按残差幅值分配并未测量“再增加一点能使真解误差降低多少”，两者不是同一量。由于真实应用没有真解，任何收益估计也必须来自可计算的代理，而不能把本实验已知解析解偷偷用于选点。

\section{参数敏感性结果的进一步解释}
宽度、权重范围和点数不是可以独立保证收敛的三个旋钮。增宽扩展名义特征空间，但若新增列高度相关，可辨识维数增长有限；增加点数可能改善覆盖，却因组平均权重改变和空间误差重新分配而导致解误差非单调；改变范围既改变特征尺度，也改变导数和矩阵谱。因此单因素扫描回答的是固定其余设置时的局部响应，不是完整的联合最优化。

在平滑直线模型上较小范围表现较好，可以解释为某些低频结构容易由缓变特征组合表示；同时小范围也可能增强列相关。这两种现象并不冲突。若只报告误差，可能忽略系数病态；若只报告条件数，又可能错过一个在指定双精度环境下确实准确的近似。本文把二者并列，避免用单个数字替代数值分析。

界面平方权重增加可加强满足界面约束的倾向，但可能牺牲内部拟合。它与增加界面点数并不等价：前者直接改变组目标的重要性，后者在组均方归一化下主要改变经验覆盖。把二者都笼统称为“加强界面损失”，会掩盖具体机制。本文所有对照均保留原权重与点数，便于单独核对。

岭参数也不是越大越稳定越好。较强收缩可以减小系数，但可能去掉真解所需方向；未经截断和经过岭的条件数必须区分数据矩阵与增广矩阵。现有扫描没有建立自动选择岭参数的方法，也没有用最终解析误差为每个种子挑参数，因此报告的是固定设置的敏感性，而非可在未知真解下直接部署的参数选择规则。

\section{近似复现中的可重复与可对齐}
可重复是指用本项目公开配置能够再次得到相近输出；与原文可对齐则要求使用相同模型、初始化、点分布、导数、权重、求解阈值和误差规则。前者并不自动推出后者。本文已经保存自身配置，但缺少作者若干实现细节，因此保留近似复现定位。

原表误差与本项目差异不能通过任意更换种子、删去较差运行或只保留较小积分值来消除。先前对初始化范围的另一解释也保留为诊断，而没有改写成作者原始方案。这种处理可能使结果表更复杂，却能使读者知道哪些设置来自原文字面，哪些来自本项目的补充假设。

解析导数与FD导数在部分设置接近，只说明在当前初始化、阈值及误差量级下，两种实现得到相近预测；不能证明FD矩阵与解析矩阵完全相同，也不能推导所有范围和步长下等价。步长扫描显示很小的差分步长可能恶化误差，符合舍入放大的机制，但其定量幅度还受线性系统敏感性影响。

花形积分误差的相对变化在最小误差组较大，提示接近极小数值时应降低对末几位的解释强度。本文不比较两方法某一末位数字来宣称优劣，也不把数值积分结果当作严格误差认证。后续若需高精度结论，应增加积分密度和高精度求值，并明确把额外成本计入。

\section{对原研究设想的回答}
原设想是利用内部与界面残差共同分配资源，以便在界面困难或高对比度时更有效地采样。现有实现确实完成了这一机制，也完成了相应基线及消融；因此研究过程仍围绕原问题展开。变化在于结果没有支持预期的稳定优势，论文据此把贡献从候选新方法收束为独立实现与系统分析。

已有证据支持三个有限判断：一是合理静态配额有时已能达到消融收益，二是内部残差下降不保证解误差和物理通量同步下降，三是自适应本身有不可忽略的完整成本。这些判断都能在表格中找到对应配置，不依赖作者主观挑选最好运行。它们也不等价于证明更精细的联合指标永远无用。

后续最有价值的改进应针对已观察的机制提出可证伪假设，例如在固定表达空间下分离配额与位置选择，或在冻结权重后使用新的监测池。若改成增长网络、更多子域或更大特征数，必须给静态基线相同新增表达能力，否则无法知道收益是否来自采样。先明确比较对象，再决定是否值得继续扩展，比不断扩大参数扫描更有解释力。

本章最终没有以“算法一定成功”作为验收标准，而是以模型与代码一致、资源计数清楚、全种子数据保留和结论可核对作为标准。尚未满足的学校形式要求及文献缺口属于终稿准备事项，不能通过把负面结果改成正面陈述来解决。


\chapter{总结与展望}
\section{主要工作与结论}
本文建立了CPU可运行、配置和原始结果可追溯的随机特征科研工程，验证Poisson五点差分、双子域解析导数及非零界面跳跃，并实现固定均匀、固定随机、PDE残差和联合指标采样。数学分析包括离散矩阵、局部尺度、加权最小二乘、标准右端扰动界及岭滤波解释。花形界面FD近似复现补充了与原文实现细节不同的诊断。

第一，合理静态基线不可缺少。当前平滑模型的联合策略没有稳定胜过均匀或随机点，固定更多界面点也可取得消融中的部分收益，不能只依据一个较差固定比例宣称方法改进。

第二，采样不能脱离表达力和目标解释。独立窄峰在当前固定特征及预算下未达到高精度，具体误差来源尚未唯一识别；联合策略在部分对照改善而未取得普遍优势。非均匀经验目标、配额和位置选择的效应需要区分，单纯增加点数的误差也不必单调。

第三，条件数、系数敏感性和物理解误差不是同一指标。尺度平衡可改善解却恶化物理通量；增宽可降误差同时增加病态；SVD截断会改变有效近似空间。固定矩阵的系数放大界有数值支持，但小函数变幅不足以证明一般稳定性。

第四，复现必须保留不一致。本文确认了原文花形数据和FD设置，却未取得足以逐值严格重现的全部细节，故明确标为近似复现。与原表更低或更高的误差都不自动构成新贡献。

\section{与最初研究目标的对应}
最低目标对应Poisson与简单界面的正确实现，已有模型一致性、单元测试和差分收敛数据支撑。标准目标对应采样对照、参数敏感性、误差与成本分析，已形成可审阅证据。拓展目标原本要求联合区域残差与界面误差的方法优于合理基线；当前证据没有达到这一条件。因此，研究方向与初衷一致，但不能把预期优势写成已实现的创新成果。

随机特征、局部网络、残差采样、SVD及界面约束均属于已有技术。本文可明确归于自身实施的工作是制造解与独立代码、共享特征的预算比较、同配额诊断、负面结果保留，以及有限维误差分解与扰动实验的组织。所给代数命题是标准工具的自含推导，不作为原创定理。尚未证明联合策略普遍优越，也未证明所有联合策略均无效；结论仅否定当前实验对强优势声明的支持。

\section{局限及后续方向}
研究仍局限于少量二维制造解、固定单隐层特征和五至十种子；正式版公式、RNN-PG全文和完整新颖性核查有缺口。实验没有建立随机自适应的全PDE误差理论，也没有完成学校官方模板确认。窄峰结果提示未来可研究更多局部子域、特征增长或算子感知采样，但这些已有相关工作\cite{ag,christoffel,schwarzpreprint}，应先明确差异，且把额外表达力和成本公平地给予对照。

后续优先处理独立监测池、尺度与重要性权重的关系、矩阵扰动及高精度抵消核查，并在新的模型上验证从已有结果形成的假设。论文的合理定位仍为既有方法的独立实现、可重复数值比较和稳定性分析；是否形成算法改进，应由更强证据决定。

% 数字顺序制审阅稿，元数据来自保存的Crossref/arXiv来源；学校细则待确认。
\begin{thebibliography}{99}
\addcontentsline{toc}{chapter}{参考文献}
\raggedright
\bibitem{dgpreprint} SUN J, DONG S, WANG F. Local randomized neural networks with discontinuous Galerkin methods for partial differential equations[EB/OL]. arXiv:2206.05577v1, 2022[2026-10-13]. \url{https://arxiv.org/abs/2206.05577v1}. 正式版：Journal of Computational and Applied Mathematics, 2024, 445:115830. DOI: \url{https://doi.org/10.1016/j.cam.2024.115830}.

\bibitem{interfacepreprint} LI Y, WANG F. Local Randomized Neural Networks Methods for Interface Problems[EB/OL]. arXiv:2308.03087v1, 2023[2026-10-13]. \url{https://arxiv.org/abs/2308.03087v1}. 正式版题名：Local randomized neural networks with finite difference methods for interface problems, Journal of Computational Physics, 2025, 529:113847. DOI: \url{https://doi.org/10.1016/j.jcp.2025.113847}. 本文算法阅读和复现依据预印本。

\bibitem{ag} DANG H, WANG F, JIANG S. Adaptive-Growth Randomized Neural Networks for PDEs: Algorithms and Numerical Analysis[EB/OL]. arXiv:2408.17225，首次2024，本文读取v3（2025-11-23）[2026-10-13]. \url{https://arxiv.org/abs/2408.17225v3}. DOI: \url{https://doi.org/10.48550/arXiv.2408.17225}.

\bibitem{schwarzpreprint} SHANG Y, HEINLEIN A, MISHRA S, WANG F. Overlapping Schwarz preconditioners for randomized neural networks with domain decomposition[EB/OL]. arXiv:2412.19207v1, 2024[2026-10-13]. \url{https://arxiv.org/abs/2412.19207v1}. 正式版：Computer Methods in Applied Mechanics and Engineering, 2025, 442:118011. DOI: \url{https://doi.org/10.1016/j.cma.2025.118011}.

\bibitem{rar} WU C, ZHU M, TAN Q, KARTHA Y, LU L. A comprehensive study of non-adaptive and residual-based adaptive sampling for physics-informed neural networks[J]. Computer Methods in Applied Mechanics and Engineering, 2023, 403:115671. DOI: \url{https://doi.org/10.1016/j.cma.2022.115671}. 本文读取公开预印本arXiv:2207.10289。

\bibitem{christoffel} LINGHU J, WANG Y. Residual-Christoffel Sampling for Random Feature Collocation of Linear PDEs[EB/OL]. arXiv:2607.13382v1, 2026[2026-10-13]. \url{https://arxiv.org/abs/2607.13382v1}.

\bibitem{pg} SHANG Y, WANG F, SUN J. Randomized neural network with Petrov--Galerkin methods for solving linear and nonlinear partial differential equations[J]. Communications in Nonlinear Science and Numerical Simulation, 2023, 127:107518. DOI: \url{https://doi.org/10.1016/j.cnsns.2023.107518}. 正式书目已核实，本文未获全文，未据此归纳方法或作者结论。

\bibitem{elm2006} HUANG G B, ZHU Q Y, SIEW C K. Extreme learning machine: Theory and applications[J]. Neurocomputing, 2006, 70(1--3):489--501. DOI: \url{https://doi.org/10.1016/j.neucom.2005.12.126}. 本次核对出版页摘要，未取得本地全文。

\bibitem{elmfbp} ANDERSON S, DOLEAN V, MOSELEY B, et al. ELM-FBPINNs: An Efficient Multilevel Random Feature Method[EB/OL]. arXiv:2409.01949v3, 2026[2026-10-16]. \url{https://arxiv.org/abs/2409.01949v3}. 首次提交2024，本文阅读第三版PDF。

\bibitem{golub} GOLUB G H, VAN LOAN C F. Matrix computations[M]. 4th ed. Johns Hopkins University Press, 2013. ISBN:9781421407944. DOI: \url{https://doi.org/10.56021/9781421407944}. 版本和年份按OpenLibrary第4版记录；本项目未取得教材全文，正文只使用并证明标准线性代数结论。

\bibitem{scipy} VIRTANEN P, GOMMERS R, OLIPHANT T E, et al. SciPy 1.0: fundamental algorithms for scientific computing in Python[J]. Nature Methods, 2020, 17:261--272. DOI: \url{https://doi.org/10.1038/s41592-019-0686-2}.
\end{thebibliography}

\appendix
\chapter{制造解、程序与复现说明}
\section{局部峰导数}
令$q(s)=s(1-s)$、$\xi=x-c_x,\eta=y-c_y,a=w^{-2}$，则
\begin{align}
h_x&=16A e^{-a(\xi^2+\eta^2)}[q'(x)-2a\xi q(x)]q(y),\\
h_{xx}&=16A e^{-a(\xi^2+\eta^2)}[-2-4a\xi q'(x)+(4a^2\xi^2-2a)q(x)]q(y).
\end{align}
$h_{yy}$交换$x,y$即可；$f_i^{\rm peak}=f_i-\beta_i\Delta h$。两侧共同$h$在值跳跃中相消，通量却因$\beta$不同而改变。

花形曲线$X=c+r(\theta)(\cos\theta,\sin\theta)$的切向为$(r'\cos\theta-r\sin\theta,r'\sin\theta+r\cos\theta)$。其外法向为切向顺时针旋转后的单位向量，分母$\sqrt{r^2+(r')^2}$。内部强迫为$-(x^2+y^2)e^{xy}$，外部强迫为$2\sin x\sin y$；外边界非零，应由外部真解计算。

\section{实验与数据对应}
项目根目录名为\path{Bachelor_Thesis_RaNN}。实验01原始数据保存在\path{results/raw}，其余保存在\path{results/runs}。本文主要子目录为：
\begin{itemize}
\item \path{exp02_interface_20261013}：界面与尺度，150次；
\item \path{exp03_sampling_20261013}：主采样和指标消融，180次；
\item \path{exp04_sensitivity_20261013}：固定比例与单因素，285次；
\item \path{exp06_peaks_20261013}：峰值与unused池，400次；
\item \path{exp07_stability_20261013}：15组谱、90组扰动；
\item \path{exp08_efficiency_20261013}：25次串行成本复测；
\item \path{exp09_reproduction_20261013}：110次初始复现/诊断假设；
\item \path{exp09_literal_20261013}：60次较贴原文字面的初始化追加。
\end{itemize}
每个run保存配置、逐次CSV、均值/样本标准差、输出系数和日志。图可由CSV独立重绘；重复run\_id拒绝覆盖。程序状态ok仅表示正常完成，不代表高精度或理论证明。论文表图由\path{experiments/build_thesis_assets.py}生成，并在\path{thesis/asset_sources.csv}记录来源。

\section{运行及编译}
Python依赖在requirements与lock文件。Windows与Linux解释器路径分别为：
\begin{verbatim}
.venv/Scripts/python.exe
.venv/bin/python
\end{verbatim}
在项目根目录运行测试与实验入口，例如：
\begin{verbatim}
python -m unittest discover -s tests -v
python -m experiments.exp09_reproduction
\end{verbatim}
新run\_id另行指定，不覆盖已经保存的实验结果。

论文推荐XeLaTeX。已查到学校本科模板的官方附件入口，但现行附件有验证码限制，尚未取得正文；有ctexbook时使用该类，否则以book、fontspec和Noto中文字体构成可编译草稿。当前排版及数字顺序参考文献组织只用于审阅，需要按学校规范调整封面、声明、页眉、目录和引用细节。所有实际编译和质量检查结果在项目文档记录，不能把编译成功理解为论文已获学校审核。

\section{本轮修订的证据生成}
在项目根目录依次运行：
\begin{verbatim}
python -m experiments.build_thesis_assets
python -m experiments.build_thesis_v3_evidence
python -m experiments.build_model_illustration
\end{verbatim}
第一项生成前期表图，第二项仅从封存采样CSV整理多指标、矩阵资源和配额轨迹，第三项从保存的解析剖面CSV生成模型示意。三者不执行新的PDE拟合；模型示意不是网络解。内容规模统计入口为\path{experiments.audit_thesis_v3}，编译及来源校验为\path{experiments.check_thesis}。

学校规范核查、材料清单和辅助范围说明分别保存在\path{docs/dlut_requirements_review.md}及\path{docs/submission_materials_checklist.md}。现有日志和数据不能替代本人的实际学习、导师指导、学校审查及签署程序。

\backmatter
\chapter*{致谢}
\addcontentsline{toc}{chapter}{致谢}
致谢内容待本人依据实际获得的指导与帮助撰写。本稿不虚构指导过程、项目资助、人员贡献或答辩经历。

\end{document}
